% Appendix fragment for cls_neutral_abl_equilibrium_awf_note_revised.tex.
% Insert after the existing appendices and before \begin{thebibliography}:
%   % Appendix fragment for cls_neutral_abl_equilibrium_awf_note_revised.tex.
% Insert after the existing appendices and before \begin{thebibliography}:
%   % Appendix fragment for cls_neutral_abl_equilibrium_awf_note_revised.tex.
% Insert after the existing appendices and before \begin{thebibliography}:
%   % Appendix fragment for cls_neutral_abl_equilibrium_awf_note_revised.tex.
% Insert after the existing appendices and before \begin{thebibliography}:
%   \input{accepted_awf_appendix}
% The main note already calls \appendix. Do not call it again here.
% Uses the main note's macros, theorem environments, labels and bibliography.

\section{AWF with CLS-compatible coefficients, a budget-determined cutoff,
and a nonconstant inner kinetic-energy profile}
\label{app:accepted-awf}

This appendix defines the accepted alternative to the constant-inner-$k$
EP-AWF construction. The objective is to preserve the standard neutral
CLS equilibrium in the outer region while allowing the kinetic energy
to vary inside the modeled wall layer. The AWF viscosity coefficients
are calibrated to that outer equilibrium, dissipation retains the
budget-determined cutoff, and the prescribed redistribution $C_k^w$
is removed. Ordinary diffusion of the transported inner $k$ then
balances the local difference between production and dissipation.
The outer CLS model, including its Yap correction, is unchanged.

The original CLS and rough-wall AWF supply the constitutive and
wall-model starting points \cite{CLS1996,SCI2006}. The coefficient
calibration, budget cutoff, retained inner transport, and curved
completion specified here are modifications of that combination.
When this variant is selected, the definitions below replace the
original wall coefficients, interior sampling choice, and
$\chi$-dependent source policy used elsewhere in the note.

\subsection{Scope, endpoint sampling, and coefficient calibration}
\label{app:accepted-awf-coefficients}

Use the high-Re outer equation convention of
\cref{sec:derivation}: molecular momentum and scalar diffusion are
omitted in the outer analytical reference, while $\nu>0$ is retained
in the AWF. This is a statement about the selected equations, not an
assertion that the same profiles solve a finite-Re outer system.
The transported inner turbulence variable is $k$; inner
dissipation is a dependent total $k$ sink. No inner CLS dissipation
transport equation is introduced.

The accepted sampling and tail-matching heights are
\begin{equation}
 \boxed{z_K=z_m=H,\qquad
 K=k_w(H)=k_o(H)>0.}
 \label{eq:accepted-sampling}
\end{equation}
The subscripts $w,o$ distinguish the modeled wall layer and outer
region. At each positive current $K$, define the counterfactual
reference scales
\begin{equation}
 \ust(K)=\sqrt{\frac{K}{\qs}},\qquad
 t_0(K)=\frac{K}{\qs},\qquad
 \Aref(K)=\frac{\ust(K)^3}{\kc}.
 \label{eq:accepted-reference-scales}
\end{equation}
These are reference scales. In a developing, forced, or separated
flow, $t_0(K)$ is not imposed as the actual wall traction.
Arguments $K$ are suppressed below when all quantities use the same
sample. A fixed inlet reference is obtained by setting $K=K_0$.

The outer constants are the roots derived in
\cref{eq:qroot,eq:kappaeq},
\begin{equation}
 \begin{gathered}
 C(\qs)\qs^2=1,\qquad
 \frac{\qs\kc^2}{1.3}
 +H_Y\!\left(\frac{\kc\qs^{3/2}}{2.5}\right)=0.48,\\
 \qs=3.190675081624288,\qquad
 \kc=0.440553905186659,\qquad
 \Cs=\qs^{-2}=0.0982278953282113.
 \end{gathered}
 \label{eq:accepted-cls-constants}
\end{equation}
Here $C$ and $H_Y$ retain the definitions of the original high-Re CLS
closure and its Yap correction. Matching the dissipation coefficient
and the asymptotic auxiliary-viscosity slope gives
\begin{equation}
 \boxed{\begin{aligned}
 c_\ell^w&=\kc\qs^{3/2}=2.510862475219797,\\
 c_\mu^w&=\Cs=\qs^{-2}=0.0982278953282113,\\
 \alpha&=c_\mu^wc_\ell^w
       =\frac{\kc}{\sqrt{\qs}}=0.246636736399424.
 \end{aligned}}
 \label{eq:accepted-coefficients}
\end{equation}
Indeed, with $a:=\alpha\sqrt K$,
\begin{equation}
 \boxed{a=\kc\ust,\qquad
 \frac{K^{3/2}}{c_\ell^w}=\Aref
 =\frac{t_0^2}{a}.}
 \label{eq:accepted-calibration-identities}
\end{equation}
The choice of $c_\ell^w$ matches the outer dissipation tail exactly.
The additional choice of $c_\mu^w$ matches the viscosity slope;
conservative interface matching by itself does not require equality
of these slopes.

Retain the original intercept correlation $b(K,h,\nu)$ in
\cref{eq:brough}, including $b=10.7$ for $h=0$. Define
\begin{equation}
 z_v=\frac{\nu b}{\sqrt K},\qquad
 \nu_{t,A}(z)=a\max(z-z_v,0),\qquad
 D_A=\nu+\nu_{t,A},\qquad
 D_{k,w}=\nu+\frac{\nu_{t,A}}{\sigma_k}.
 \label{eq:accepted-viscosities}
\end{equation}
The CLS choice $\sigma_k=1$ gives $D_{k,w}=D_A$.
These auxiliary coefficients depend on sampled $K$, while the
transported field $k_w(z)$ can vary. Substituting $k_w(z)$ for $K$
in the viscosity or dissipation law would define a different model.

\subsection{Momentum, production, and the outer velocity anchor}
\label{app:accepted-awf-momentum}

The analytical form of the AWF momentum reconstruction is retained,
with all integrals evaluated using the calibrated $D_A$. Put
$L_h(z)=\max(z-h,0)$ and
\begin{equation}
 J_0(z)=\int_0^z\frac{\dd s}{D_A(s)},\qquad
 J_1(z)=\int_0^z\frac{L_h(s)}{D_A(s)}\,\dd s,\qquad
 J_2(z)=\int_0^z\frac{\dd s}{D_A(s)^2}.
 \label{eq:accepted-J}
\end{equation}
For scalar tangential forcing $F$, the flat quasi-steady reduction is
\begin{equation}
 \begin{aligned}
 t_A&=\frac{U_H-FJ_1(H)}{J_0(H)},\qquad
 U_A(z)=t_AJ_0(z)+FJ_1(z),\\
 P_A(z)&=\nu_{t,A}(z)
       \left[\frac{t_A+FL_h(z)}{D_A(z)}\right]^2.
 \end{aligned}
 \label{eq:accepted-forced-awf}
\end{equation}
Thus no slip, the original forcing cutoff at $h$, and the explicit
pressure-gradient response retain their analytical form. Their
numerical values change with the new viscosity coefficients.
The vector form follows \cref{eq:3d}.

At the counterfactual zero-forcing reference, choose
\begin{equation}
 U_{H,0}=t_0J_0(H),\qquad
 U_{w,0}(z)=t_0J_0(z),\qquad
 p_{A,0}(z)=t_0^2\frac{\nu_{t,A}(z)}{D_A(z)^2}.
 \label{eq:accepted-reference-production}
\end{equation}
The resulting integrated reference production is
\begin{equation}
 \boxed{\mathcal P_0:=\int_0^H p_{A,0}(z)\,\dd z
 =t_0^2[J_0(H)-\nu J_2(H)]=HP_0.}
 \label{eq:accepted-production-budget}
\end{equation}
$P_0$ remains the height mean used in the main note; $p_{A,0}$ is
pointwise reference production. Neither is the production of a
general current terrain flow.

For $z\geq H$, preserve the standard outer CLS branch
\begin{equation}
 \boxed{\begin{aligned}
 U_{o,0}(z)&=t_0J_0(H)+\frac{\ust}{\kc}\ln\frac zH,\\
 k_{o,0}(z)&=K,\qquad
 \eps_{o,0}(z)=\frac{\Aref}{z},\qquad
 \nu_{t,o,0}(z)=\kc\ust z.
 \end{aligned}}
 \label{eq:accepted-outer-reference}
\end{equation}
The outer velocity offset has been used to satisfy the AWF anchor.
Its aerodynamic intercept is therefore
\begin{equation}
 z_{0,c}=H\exp[-\kc\ust J_0(H)].
 \label{eq:accepted-aerodynamic-intercept}
\end{equation}
Retaining the same sand-grain $h$ while changing the AWF coefficients
generally changes this intercept. A separately prescribed $z_{0,c}$
requires an additional momentum/intercept calibration.

\subsection{Budget-determined dissipation and admissibility}
\label{app:accepted-awf-cutoff}

Use the continuous positive dissipation distribution
\begin{equation}
 \boxed{E_w(z)=\ECw(z)=\frac{\Aref}{\max(z,z_c)},
 \qquad C_k^w(z)=0.}
 \label{eq:accepted-dissipation}
\end{equation}
The required tail begins no later than $H$, so $0<z_c\leq H$.
Its integral is
\begin{equation}
 \int_0^H\ECw(z)\,\dd z
 =\Aref\left(1+\ln\frac H{z_c}\right).
 \label{eq:accepted-dissipation-integral}
\end{equation}
The standard outer reference has zero $k$ gradient. Together with
zero modeled ground $k$ flux, conservative matching therefore
requires
\begin{equation}
 \boxed{\int_0^H\ECw\,\dd z=\mathcal P_0.}
 \label{eq:accepted-budget-condition}
\end{equation}
Combining the preceding identities gives
\begin{equation}
 \boxed{z_c=H\exp(1-\Lambda),\qquad
 \Lambda=\frac{\mathcal P_0}{\Aref}
 =\frac{t_0^2[J_0(H)-\nu J_2(H)]}{\Aref}.}
 \label{eq:accepted-cutoff}
\end{equation}
The left side of \cref{eq:accepted-dissipation-integral} decreases
strictly with $z_c$, tends to infinity as $z_c\downarrow0$, and equals
$\Aref$ at $z_c=H$. Consequently,
\begin{equation}
 \boxed{0<z_c\leq H\quad\Longleftrightarrow\quad\Lambda\geq1,}
 \label{eq:accepted-cutoff-gate}
\end{equation}
and the admissible cutoff is unique. This is the $z_m=H$
specialization of \cref{eq:hp-gate}. If the gate fails, this
monotone cutoff family cannot meet both the budget and outer trace.
Clipping $z_c$ to $H$ would not preserve the budget.

The original molecular cutoff
$z_{\eps,A}=2c_\ell^w\nu/\sqrt K$ is replaced by
\cref{eq:accepted-cutoff}. The plateau is now $\Aref/z_c$.
It cannot simultaneously remain $2\nu K/z_c^2$: imposing both
plateau formulas would force the original molecular cutoff again.
Thus the cutoff change is an explicit change of the inner
dissipation model.

\subsection{Derivation of the nonconstant inner reference}
\label{app:accepted-awf-inner-k}

Retain the inner transport equation with its physical production
and prescribed dissipation,
\begin{equation}
 \partial_t k_w+\bm U_w\cdot\nabla k_w
 =\nabla\cdot(D_{k,w}\nabla k_w)+P_A-\ECw.
 \label{eq:accepted-k-transport}
\end{equation}
For the steady, parallel flat reference this reduces to
\begin{equation}
 (D_{k,w}k_{w,0}')'=\ECw-p_{A,0},\qquad
 D_{k,w}(0)k_{w,0}'(0)=0,\qquad k_{w,0}(H)=K.
 \label{eq:accepted-k-bvp}
\end{equation}
Define the cumulative source defect
\begin{equation}
 Q_0(z)=\int_0^z[\ECw(s)-p_{A,0}(s)]\,\dd s.
 \label{eq:accepted-Q}
\end{equation}
Two integrations give the unique solution for the prescribed
positive diffusivity,
\begin{equation}
 \boxed{D_{k,w}(z)k_{w,0}'(z)=Q_0(z),\qquad
 k_{w,0}(z)=K-\int_z^H\frac{Q_0(s)}{D_{k,w}(s)}\,\dd s.}
 \label{eq:accepted-k-profile}
\end{equation}
The quantity $Q_0$ is the gradient flux inside the diffusion
operator; the physical upward diffusive flux is $-Q_0$.
The cutoff budget gives
\begin{equation}
 Q_0(0)=Q_0(H)=0,\qquad k_{w,0}'(H)=0.
 \label{eq:accepted-endpoint-fluxes}
\end{equation}
Conversely, these two zero endpoint fluxes imply
\cref{eq:accepted-budget-condition} by integration of the $k$
equation. Allowing a variable $k_w$ supplies local diffusion; it
does not remove the integrated compatibility requirement.

No additional value of $k_w(0)$ is prescribed. It follows from
\cref{eq:accepted-k-profile}. This is a modeled rough-wall layer,
not a viscous-sublayer resolution with an independently imposed
$k(0)=0$. The formula defines a stationary reference and an
initialization; during a general flow, $k_w$ is evolved by
\cref{eq:accepted-k-transport}, rather than reset to this profile.

For every sufficiently regular test function $\phi$, integration
by parts yields
\begin{equation}
 \boxed{-\int_0^H D_{k,w}k_{w,0}'\phi'\,\dd z
 +\int_0^H(p_{A,0}-\ECw)\phi\,\dd z=0.}
 \label{eq:accepted-weak-equilibrium}
\end{equation}
Thus the nonconstant reference balances every continuum test
moment. The higher-mode obstruction in \cref{sec:hp-source}
assumed constant inner $k$, for which the diffusion term vanished.
Finite-dimensional projection and numerical differentiation can
still leave a discrete residual; continuum balance alone does not
establish exact spectral/$hp$ preservation.

\subsection{Closed formulas and positivity on the rough branch}
\label{app:accepted-awf-rough}

Suppose $z_v\leq0$ and define
\begin{equation}
 d=\nu-az_v\geq\nu,\qquad D_A(z)=d+az.
 \label{eq:accepted-rough-D}
\end{equation}
Then
\begin{equation}
 \begin{aligned}
 J_0(z)&=\frac1a\ln\frac{d+az}{d},&
 J_2(z)&=\frac1a\left(\frac1d-\frac1{d+az}\right),\\
 p_{A,0}(z)&=t_0^2\frac{d+az-\nu}{(d+az)^2}.&&
 \end{aligned}
 \label{eq:accepted-rough-integrals}
\end{equation}
Using $\Aref=t_0^2/a$, the cumulative production is
\begin{equation}
 \int_0^z p_{A,0}(s)\,\dd s
 =\Aref\left[
 \ln\frac{d+az}{d}
 -\nu\left(\frac1d-\frac1{d+az}\right)\right].
 \label{eq:accepted-rough-production-integral}
\end{equation}
In particular, with $\zeta_H=aH/d$ and $\beta=\nu/d$,
\begin{equation}
 \Lambda=\ln(1+\zeta_H)
          -\frac{\beta\zeta_H}{1+\zeta_H}.
 \label{eq:accepted-rough-Lambda}
\end{equation}
The cumulative dissipation is $\Aref R_c(z)$, where
\begin{equation}
 R_c(z)=\begin{cases}
 z/z_c,&0\leq z\leq z_c,\\
 1+\ln(z/z_c),&z_c\leq z\leq H.
 \end{cases}
 \label{eq:accepted-Rc}
\end{equation}
Hence an elementary expression for the flux is
\begin{equation}
 Q_0(z)=\Aref\left[
 R_c(z)-\ln\frac{d+az}{d}
 +\nu\left(\frac1d-\frac1{d+az}\right)\right].
 \label{eq:accepted-rough-Q}
\end{equation}
Together with \cref{eq:accepted-k-profile}, this provides an exact
analytical representation of the inner reference.

\paragraph{A closed primitive for the kinetic energy.}
For the selected $\sigma_k=1$, set
$\zeta=az/d$, $\zeta_c=az_c/d$, and define
\begin{align}
 \mathcal F_< (\zeta)
 &=\frac{\zeta-\ln(1+\zeta)}{\zeta_c}
   -\frac12\ln^2(1+\zeta)
   +\beta\left[\ln(1+\zeta)+\frac1{1+\zeta}-1\right],
 \label{eq:accepted-primitive-lower}\\
 \mathcal G(\zeta)
 &=\left[1+\ln\frac{\zeta}{\zeta_c}+\beta\right]\ln(1+\zeta)
   +\operatorname{Li}_2(-\zeta)
   -\frac12\ln^2(1+\zeta)+\frac{\beta}{1+\zeta},
 \quad \zeta>0.
 \label{eq:accepted-primitive-upper}
\end{align}
The real dilogarithm is defined here by
$\operatorname{Li}_2(-\zeta)=-\int_0^\zeta\ln(1+s)\,\dd s/s$.
Join the primitives continuously:
\begin{equation}
 \mathcal F(\zeta)=
 \begin{cases}
 \mathcal F_< (\zeta),&0\leq\zeta\leq\zeta_c,\\
 \mathcal F_< (\zeta_c)+\mathcal G(\zeta)-\mathcal G(\zeta_c),
 &\zeta_c\leq\zeta\leq\zeta_H.
 \end{cases}
 \label{eq:accepted-primitive-joined}
\end{equation}
Direct differentiation of these expressions gives
$\mathcal F'(\zeta)=Q_0(d\zeta/a)/[\Aref(1+\zeta)]$.
Therefore the closed profile is
\begin{equation}
 \boxed{k_{w,0}(z)=K+\frac{\Aref}{a}
 \left[\mathcal F\!\left(\frac{az}{d}\right)
       -\mathcal F(\zeta_H)\right].}
 \label{eq:accepted-k-closed}
\end{equation}
The integral form \cref{eq:accepted-k-profile} remains applicable
to other viscosity segments and positive diffusivities.

\paragraph{The sign of the local imbalance.}
For every $z\geq z_c>0$ on this rough branch,
\begin{equation}
 \boxed{\ECw(z)-p_{A,0}(z)
 =\frac{\Aref[(d+\nu)az+d^2]}{z(d+az)^2}>0.}
 \label{eq:accepted-tail-defect}
\end{equation}
This strict inequality concerns the $1/z$ tail. It does not imply
positive imbalance throughout the plateau. In fact, the budget
requires a compensating region of negative imbalance.

\begin{proposition}[A positive rough-wall reference]
Assume $z_v\leq0$, $d\geq2\nu$, and $\Lambda\geq1$.
Then $\ECw-p_{A,0}$ has exactly one zero in $(0,z_c)$,
changing from negative to positive. For every positive $D_{k,w}$,
\begin{equation}
 Q_0(z)<0\quad(0<z<H),\qquad
 k_{w,0}'(z)<0\quad(0<z<H),\qquad
 k_{w,0}(z)>K>0\quad(0\leq z<H).
 \label{eq:accepted-positivity}
\end{equation}
\end{proposition}
\begin{proof}
The reference production satisfies
\[
 p_{A,0}'(z)=
 \frac{a t_0^2[2\nu-(d+az)]}{(d+az)^3}\leq0,
\]
and is strictly decreasing for $z>0$. On the dissipation plateau,
$\Aref/z_c-p_{A,0}$ is therefore strictly increasing and is
positive at $z_c$ by \cref{eq:accepted-tail-defect}.
Its initial value must be negative, since otherwise the total
integral could not vanish. It crosses zero exactly once, after
which the defect remains positive up to $H$. Its cumulative
integral first decreases and then increases to $Q_0(H)=0$, so
$Q_0<0$ in the interior. The profile and monotonicity statements
follow from \cref{eq:accepted-k-profile}.
\end{proof}
The condition $d\geq2\nu$ is sufficient, not necessary.
For other intercepts, including a smooth-wall viscosity segment,
the budget construction remains defined when its gate holds, but
positivity must be assessed from the resulting profile.

The rough-branch velocity anchor also gives the useful identities
\begin{equation}
 \boxed{z_{0,c}=\frac{dH}{aH+d},\qquad
 z_c=\ee z_{0,c}
 \exp\!\left[\nu\left(\frac1d-\frac1{aH+d}\right)\right].}
 \label{eq:accepted-roughness-cutoff}
\end{equation}
Even with the calibrated slope, $z_{0,c}$ retains a finite-height
dependence. For $aH\gg d$, it tends to $d/a$.

\subsection{Conservative matching and normal pressure balance}
\label{app:accepted-awf-matching}

The inner momentum reference and the outer profile in
\cref{eq:accepted-outer-reference} satisfy
\begin{equation}
 \begin{gathered}
 U_{w,0}(0)=0,\qquad U_{w,0}(H)=U_{o,0}(H),\\
 D_A(H)U_{w,0}'(H)
 =\nu_{t,o,0}(H)U_{o,0}'(H)=t_0.
 \end{gathered}
 \label{eq:accepted-momentum-match}
\end{equation}
The calibrated coefficients match the viscosity slope, not the
finite-height viscosities: on the rough branch,
$D_A(H)=aH+d$ while $\nu_{t,o,0}(H)=aH$.
The velocity derivatives therefore need not coincide.

The turbulence matching conditions are
\begin{equation}
 \boxed{\begin{aligned}
 k_{w,0}(H)&=k_{o,0}(H)=K,\\
 D_{k,w}(H)k_{w,0}'(H)
 &=D_{k,o}(H)k_{o,0}'(H)=0,\\
 \ECw(H)&=\eps_{o,0}(H)=\frac{\Aref}{H}.
 \end{aligned}}
 \label{eq:accepted-turbulence-match}
\end{equation}
Here $D_{k,o}=\nu_t^{CLS}/\sigma_k$ on the declared outer branch.
Inner dissipation is dependent, so the outer dissipation equation
receives its single trace. No additional dissipation-flux condition
is imposed. In the limiting case $z_c=H$, even the one-sided
dissipation derivatives differ, while the required trace still
matches.

With the full-strain inner stress completion of
\cref{eq:curve-inner-stress}, the parallel inner normal stress is
$R_{zz,w}=2k_w/3$. Its normal momentum balance therefore requires
\begin{equation}
 \frac{\dd}{\dd z}\left(
 \frac{p_w}{\rho}+gz+\frac23 k_{w,0}\right)=0.
 \label{eq:accepted-inner-pressure}
\end{equation}
Mechanical pressure inside this layer is generally not hydrostatic.
Since $R_{zz,o}=0.52K$ at the outer reference, continuity of total
normal traction gives
\begin{equation}
 \boxed{\frac{p_w(H)-p_o(H)}{\rho}
 =\left(0.52-\frac23\right)K=-\frac{11}{75}K.}
 \label{eq:accepted-pressure-jump}
\end{equation}
The interface offset is unchanged from the full-strain EP-AWF
completion because the matching value is still $K$; the interior
pressure distribution changes with $k_{w,0}(z)$.
When a reduced pressure is stored, recover the mechanical pressure
consistently before evaluating the AWF pressure forcing.

The flat reference also retains the integrated mechanical-energy
identity
\begin{equation}
 t_0 U_{H,0}
 =\nu\int_0^H[U_{w,0}'(z)]^2\,\dd z
  +\int_0^H p_{A,0}(z)\,\dd z
 =\nu t_0^2J_2(H)+\int_0^H\ECw(z)\,\dd z.
 \label{eq:accepted-energy}
\end{equation}
The nonzero internal $k$ flux redistributes energy within the layer;
it contributes no net flux at the two reference endpoints.

\subsection{Distinction from EP-AWF with unit equilibrium weight}
\label{app:accepted-awf-versus-ep}

To isolate the transport difference, compare the two constructions
at identical $K,h,H,\nu$, viscosity coefficients, and budget cutoff.
At $\chi=1$, the EP-AWF has
\begin{equation}
 E_w^{EP}=\ECw,\qquad
 C_{k,EP}^w=\ECw-p_{A,0},\qquad
 S_{k,w}^{EP}=P_A-p_{A,0}.
 \label{eq:accepted-ep-comparison}
\end{equation}
At the flat reference, $P_A=p_{A,0}$ and this source vanishes
pointwise, permitting constant inner $k=K$. The accepted alternative
instead uses
\begin{equation}
 E_w=\ECw,\qquad C_k^w=0,\qquad
 S_{k,w}=P_A-\ECw.
 \label{eq:accepted-source-comparison}
\end{equation}
Its reference source is generally nonzero and is balanced by
$\partial_z(D_{k,w}k_{w,0}')$.

Both sources have the same layer integral,
\begin{equation}
 \int_0^H S_{k,w}\,\dd z
 =\int_0^H S_{k,w}^{EP}\,\dd z
 =\int_0^H(P_A-p_{A,0})\,\dd z,
 \label{eq:accepted-equal-integrated-sources}
\end{equation}
but their local balances and inner kinetic-energy fields differ.
On the reference, EP-AWF carries the physical redistribution flux
$-Q_0$ through its prescribed correction; the accepted model carries
the same flux through ordinary diffusion of its nonconstant $k$.
This reference identity does not make their evolving dynamics
equivalent. In addition, comparison with an EP-AWF that retains the
original coefficients must account separately for the changed
momentum reconstruction.

Setting the old switch to $\chi=0$ does not select the present
variant: it restores the original-cutoff dissipation $\EAw$ through
\cref{eq:blend}. The accepted model retains $\ECw$ while removing
$C_k^w$. Consequently the spatial fade to the original AWF in
\cref{sec:transition} is not part of this variant.

\subsection{A specified extension to curved terrain}
\label{app:accepted-awf-curved}

The accepted scalar policy can be used with the conservative curved
patch of \cref{sec:curved}. To make that extension definite, adopt
the note's metric-weighted reference budget, together with
$r_K=r_m=H$, $K=k|_{\Gamma_H}$, the calibrated viscosity, and
$C_k^w=0$. This is an explicit extension of the flat model; the
flat matching theorem does not by itself select a unique curved
calibration or establish an equilibrium over a hill.

Retain the plus curvature convention
\begin{equation}
 \bm x(\bm\xi,r)=\bm X(\bm\xi)+r\bm n(\bm\xi),\qquad
 \mathsf B=\nabla_\Gamma\bm n,\qquad
 \mathsf A=I_T+r\mathsf B,\qquad J=\det_T\mathsf A.
 \label{eq:accepted-curved-map}
\end{equation}
Here $r$ is physical normal distance and $\dd V=J\,\dd S_w\,\dd r$.
All principal stretches $1+r\kappa_i$ must stay positive, and
the normal map must be globally one-to-one. Positive local $J$
alone does not guarantee unique normal projection.

\paragraph{Reference cutoff in the physical volume measure.}
Define
\begin{equation}
 \boxed{\begin{aligned}
 E_w(r)&=\ECw(r)=\frac{\Aref(K)}{\max(r,r_c)},\\
 \int_0^H J(r)\ECw(r)\,\dd r
 &=\mathcal P_J:=\int_0^H J(r)p_{A,0}(r;K)\,\dd r,\\
 p_{A,0}(r;K)&=t_0(K)^2
              \frac{\nu_{t,A}(r;K)}{D_A(r;K)^2}.
 \end{aligned}}
 \label{eq:accepted-curved-budget}
\end{equation}
The reference kernel remains counterfactual flat, parallel,
zero-forcing production at the current $K$. It is not replaced
by actual terrain production.
For $\mathcal B(c)=\int_0^H J(r)\Aref/\max(r,c)\,\dd r$,
\begin{equation}
 \mathcal B'(c)=-\frac{\Aref}{c^2}\int_0^cJ(r)\,\dd r<0,
 \qquad \lim_{c\downarrow0}\mathcal B(c)=+\infty.
 \label{eq:accepted-curved-monotonicity}
\end{equation}
Therefore a unique $0<r_c\leq H$ exists exactly when
\begin{equation}
 \boxed{\mathcal P_J\geq
 \frac{\Aref}{H}\int_0^HJ(r)\,\dd r.}
 \label{eq:accepted-curved-gate}
\end{equation}
At a fixed footpoint on a three-dimensional surface,
$J=(1+\kappa_1r)(1+\kappa_2r)$. Direct integration gives
\begin{equation}
 \begin{aligned}
 \frac{\mathcal B(c)}{\Aref}
 &=1+\ln\frac Hc
 +(\kappa_1+\kappa_2)\left(H-\frac c2\right)\\
 &\quad+\kappa_1\kappa_2
       \left(\frac{H^2}{2}-\frac{c^2}{6}\right).
 \end{aligned}
 \label{eq:accepted-curved-budget-explicit}
\end{equation}
The curved cutoff thus requires one monotone scalar root solve.
The exponential formula is recovered when $J=1$.

\paragraph{Actual inner flow and transport.}
Write $\bm U=\bm u_t+w\bm n$ and
$\bm v=\mathsf A^{-1}\bm u_t$. The inner equation is
\begin{equation}
 \boxed{\begin{aligned}
 J\partial_t k
 &+\partial_r\{J(wk-D_{k,w}\partial_rk)\}\\
 &+\operatorname{div}_\Gamma\left\{
 J[\bm v k-D_{k,w}\mathsf A^{-1}\mathsf A^{-T}
                         \nabla_\Gamma k]\right\}
 =J(P_A-\ECw).
 \end{aligned}}
 \label{eq:accepted-curved-k}
\end{equation}
Use the full continuity and momentum equations
\cref{eq:curve-continuity,eq:curve-momentum}, the stated roughness
force in \cref{eq:curve-rough-force}, and the inner completion
\begin{equation}
 \begin{gathered}
 \R_A=\tfrac23k\I-\nu_{t,A}\Sm_w,\qquad
 \mathsf T_w=-\left(\frac{p_w}{\rho}+\tfrac23k\right)\I
              +D_A\Sm_w,\\
 P_A=-\R_A:\nabla\bm U
     =\tfrac12\nu_{t,A}\Sm_w:\Sm_w
 \quad\text{for }\nabla\cdot\bm U=0.
 \end{gathered}
 \label{eq:accepted-curved-stress}
\end{equation}
The isotropic stress uses local transported $k$, not sampled $K$.
For a weakly incompressible discrete field, the production
convention discussed after \cref{eq:curve-inner-stress} retains
the corresponding $-(2k/3)\nabla\cdot\bm U$ contribution when
required by the discrete stress-work balance.

The base-normal tangential traction contains curvature and normal
motion,
\begin{equation}
 \bm t_A=D_A\left[
 \partial_r\bm u_t+\mathsf A^{-T}\nabla_\Gamma w
 -\mathsf B\mathsf A^{-1}\bm u_t\right].
 \label{eq:accepted-curved-traction}
\end{equation}
Normal momentum must likewise retain curvature acceleration and
pressure balance. Merely weighting independent flat momentum
columns by $J$ does not supply these terms. General curved and
developing flow uses the coupled patch solve of \cref{sec:curved}.

The reference budget preserves the actual non-equilibrium source:
\begin{equation}
 \int_0^H J(P_A-\ECw)\,\dd r
 =\int_0^H J(P_A-p_{A,0})\,\dd r.
 \label{eq:accepted-curved-actual-budget}
\end{equation}
This quantity need not vanish; storage, tangential transport, and
interface flux balance it. Determining the cutoff from actual
$P_A$ instead of $p_{A,0}$ would impose a different energy budget.

An auxiliary frozen normal-column construction can still be written,
\begin{equation}
 \begin{aligned}
 Q_J(r)&=\int_0^rJ(s)[\ECw(s)-p_{A,0}(s)]\,\dd s,\\
 k_{\rm col}(r)&=K-\int_r^H
            \frac{Q_J(s)}{J(s)D_{k,w}(s)}\,\dd s.
 \end{aligned}
 \label{eq:accepted-curved-column}
\end{equation}
It balances only the reference source and normal diffusion.
It is not an exact terrain-flow equilibrium, since actual
production, advection, and tangential diffusion generally differ.
Its zero base-normal endpoint flux also does not establish zero
flux through an arbitrary sloping matching surface.

\paragraph{Interface conditions.}
Retain the common velocity, total traction, and mass flux in
\cref{eq:curve-momentum-interface,eq:curve-massflux}, together with
the common $k$ trace and advective--diffusive flux in
\cref{eq:curve-k-interface}. For varying $H(\bm\xi)$, these use
\begin{equation}
 \bm N_H=\frac{\bm n-\mathsf A_H^{-T}\nabla_\Gamma H}
 {\sqrt{1+|\mathsf A_H^{-T}\nabla_\Gamma H|^2}}.
 \label{eq:accepted-curved-interface-normal}
\end{equation}
Surface derivatives in the mapped equations are taken at fixed
$r$ before restriction to $H$; in particular,
$\nabla_\Gamma K=(\nabla_\Gamma k)|_H
+(\partial_rk)|_H\nabla_\Gamma H$.
The outer dissipation equation receives only
\begin{equation}
 \eps_{o,H}=\ECw(H)=\frac{\Aref(K)}{H}.
 \label{eq:accepted-curved-epsilon-trace}
\end{equation}
The full normal-traction condition replaces the flat pressure
offset over terrain. There is no independently solved inner
dissipation equation or additional prescribed dissipation flux.

\subsection{Illustrative reference and implementation implications}
\label{app:accepted-awf-example}

For a flat reference with
\begin{equation}
 \ust=0.4\units{m\,s^{-1}},\quad
 \nu=1.5\times10^{-5}\units{m^2\,s^{-1}},\quad
 h=0.1\units{m},\quad H=2\units{m},\quad z_K=z_m=H,
 \label{eq:accepted-example-inputs}
\end{equation}
the unchanged roughness correlation and calibrated coefficients give
the values in \cref{tab:accepted-reference}. They evaluate the
analytical construction; they are not a solver-validation result.

\begin{table}[htbp]
\centering\small
\begin{tabular}{@{}lll@{}}
\toprule
Quantity & Value & Units\\
\midrule
$K$ & $0.5105080131$ & $\mathrm{m^2\,s^{-2}}$\\
$z_v$ & $-0.0052721312$ & $\mathrm m$\\
$a$ & $0.1762215621$ & $\mathrm{m\,s^{-1}}$\\
$d$ & $0.0009440632$ & $\mathrm{m^2\,s^{-1}}$\\
$\mathcal P_0$ & $0.8584509755$ & $\mathrm{m^3\,s^{-3}}$\\
$z_c$ & $0.0147555946$ & $\mathrm m$\\
$U_{H,0}$ & $5.3797062304$ & $\mathrm{m\,s^{-1}}$\\
$z_{0,c}$ & $0.0053429396$ & $\mathrm m$\\
$k_{w,0}(0)$ & $1.0229104395$ & $\mathrm{m^2\,s^{-2}}$\\
$k_{w,0}(0.1\units m)$ & $0.5447982810$ & $\mathrm{m^2\,s^{-2}}$\\
$k_{w,0}(1\units m)$ & $0.5111909311$ & $\mathrm{m^2\,s^{-2}}$\\
\bottomrule
\end{tabular}
\caption{Accepted flat-wall reference with endpoint sampling.
The outer kinetic energy remains $K$, while ordinary diffusion
supports the larger inner values.}
\label{tab:accepted-reference}
\end{table}

Here $d>2\nu$, so the positivity proposition applies. In particular,
$\ECw(0)\simeq9.84519\units{m^2\,s^{-3}}$ is smaller than
$p_{A,0}(0)\simeq26.68598\units{m^2\,s^{-3}}$.
The defect changes sign near $z=0.00931273\units m$ and remains
positive thereafter, while its full-layer integral is zero.
These signs refer to the accepted budget cutoff; retaining the
original molecular cutoff would give a different imbalance.

Implementation should evaluate $K$ from the common matching trace,
construct $D_A$ with the calibrated coefficients, and determine
the cutoff from reference production at that same $K$.
Actual production is used in \cref{eq:accepted-k-transport} or
\cref{eq:accepted-curved-k}. Initialize a flat reference with
\cref{eq:accepted-k-profile}, retain zero modeled ground $k$ flux,
and couple the matching value and conservative flux to outer CLS.
Do not reintroduce the EP redistribution while using this profile.
Quadrature should resolve each active viscosity, forcing, and
dissipation breakpoint, and the full assembled reference residual
must be checked with the actual interface and pressure operators.

The derivation establishes continuum compatibility with the
standard outer CLS equilibrium on flat, undisturbed terrain and
specifies a conservative curved extension. It does not require
the outer flow over a hill to remain logarithmic. The positive
flat reference is also not a positivity or stability theorem for
the evolving equations: a sink controlled by sampled $K$ can
remain positive as local $k_w$ decreases, and the constitutive
limitations of the unchanged outer CLS model identified in
\cref{sec:stability} remain. Any active robustness modification
must be assessed against this reference and declared separately.

% The main note already calls \appendix. Do not call it again here.
% Uses the main note's macros, theorem environments, labels and bibliography.

\section{AWF with CLS-compatible coefficients, a budget-determined cutoff,
and a nonconstant inner kinetic-energy profile}
\label{app:accepted-awf}

This appendix defines the accepted alternative to the constant-inner-$k$
EP-AWF construction. The objective is to preserve the standard neutral
CLS equilibrium in the outer region while allowing the kinetic energy
to vary inside the modeled wall layer. The AWF viscosity coefficients
are calibrated to that outer equilibrium, dissipation retains the
budget-determined cutoff, and the prescribed redistribution $C_k^w$
is removed. Ordinary diffusion of the transported inner $k$ then
balances the local difference between production and dissipation.
The outer CLS model, including its Yap correction, is unchanged.

The original CLS and rough-wall AWF supply the constitutive and
wall-model starting points \cite{CLS1996,SCI2006}. The coefficient
calibration, budget cutoff, retained inner transport, and curved
completion specified here are modifications of that combination.
When this variant is selected, the definitions below replace the
original wall coefficients, interior sampling choice, and
$\chi$-dependent source policy used elsewhere in the note.

\subsection{Scope, endpoint sampling, and coefficient calibration}
\label{app:accepted-awf-coefficients}

Use the high-Re outer equation convention of
\cref{sec:derivation}: molecular momentum and scalar diffusion are
omitted in the outer analytical reference, while $\nu>0$ is retained
in the AWF. This is a statement about the selected equations, not an
assertion that the same profiles solve a finite-Re outer system.
The transported inner turbulence variable is $k$; inner
dissipation is a dependent total $k$ sink. No inner CLS dissipation
transport equation is introduced.

The accepted sampling and tail-matching heights are
\begin{equation}
 \boxed{z_K=z_m=H,\qquad
 K=k_w(H)=k_o(H)>0.}
 \label{eq:accepted-sampling}
\end{equation}
The subscripts $w,o$ distinguish the modeled wall layer and outer
region. At each positive current $K$, define the counterfactual
reference scales
\begin{equation}
 \ust(K)=\sqrt{\frac{K}{\qs}},\qquad
 t_0(K)=\frac{K}{\qs},\qquad
 \Aref(K)=\frac{\ust(K)^3}{\kc}.
 \label{eq:accepted-reference-scales}
\end{equation}
These are reference scales. In a developing, forced, or separated
flow, $t_0(K)$ is not imposed as the actual wall traction.
Arguments $K$ are suppressed below when all quantities use the same
sample. A fixed inlet reference is obtained by setting $K=K_0$.

The outer constants are the roots derived in
\cref{eq:qroot,eq:kappaeq},
\begin{equation}
 \begin{gathered}
 C(\qs)\qs^2=1,\qquad
 \frac{\qs\kc^2}{1.3}
 +H_Y\!\left(\frac{\kc\qs^{3/2}}{2.5}\right)=0.48,\\
 \qs=3.190675081624288,\qquad
 \kc=0.440553905186659,\qquad
 \Cs=\qs^{-2}=0.0982278953282113.
 \end{gathered}
 \label{eq:accepted-cls-constants}
\end{equation}
Here $C$ and $H_Y$ retain the definitions of the original high-Re CLS
closure and its Yap correction. Matching the dissipation coefficient
and the asymptotic auxiliary-viscosity slope gives
\begin{equation}
 \boxed{\begin{aligned}
 c_\ell^w&=\kc\qs^{3/2}=2.510862475219797,\\
 c_\mu^w&=\Cs=\qs^{-2}=0.0982278953282113,\\
 \alpha&=c_\mu^wc_\ell^w
       =\frac{\kc}{\sqrt{\qs}}=0.246636736399424.
 \end{aligned}}
 \label{eq:accepted-coefficients}
\end{equation}
Indeed, with $a:=\alpha\sqrt K$,
\begin{equation}
 \boxed{a=\kc\ust,\qquad
 \frac{K^{3/2}}{c_\ell^w}=\Aref
 =\frac{t_0^2}{a}.}
 \label{eq:accepted-calibration-identities}
\end{equation}
The choice of $c_\ell^w$ matches the outer dissipation tail exactly.
The additional choice of $c_\mu^w$ matches the viscosity slope;
conservative interface matching by itself does not require equality
of these slopes.

Retain the original intercept correlation $b(K,h,\nu)$ in
\cref{eq:brough}, including $b=10.7$ for $h=0$. Define
\begin{equation}
 z_v=\frac{\nu b}{\sqrt K},\qquad
 \nu_{t,A}(z)=a\max(z-z_v,0),\qquad
 D_A=\nu+\nu_{t,A},\qquad
 D_{k,w}=\nu+\frac{\nu_{t,A}}{\sigma_k}.
 \label{eq:accepted-viscosities}
\end{equation}
The CLS choice $\sigma_k=1$ gives $D_{k,w}=D_A$.
These auxiliary coefficients depend on sampled $K$, while the
transported field $k_w(z)$ can vary. Substituting $k_w(z)$ for $K$
in the viscosity or dissipation law would define a different model.

\subsection{Momentum, production, and the outer velocity anchor}
\label{app:accepted-awf-momentum}

The analytical form of the AWF momentum reconstruction is retained,
with all integrals evaluated using the calibrated $D_A$. Put
$L_h(z)=\max(z-h,0)$ and
\begin{equation}
 J_0(z)=\int_0^z\frac{\dd s}{D_A(s)},\qquad
 J_1(z)=\int_0^z\frac{L_h(s)}{D_A(s)}\,\dd s,\qquad
 J_2(z)=\int_0^z\frac{\dd s}{D_A(s)^2}.
 \label{eq:accepted-J}
\end{equation}
For scalar tangential forcing $F$, the flat quasi-steady reduction is
\begin{equation}
 \begin{aligned}
 t_A&=\frac{U_H-FJ_1(H)}{J_0(H)},\qquad
 U_A(z)=t_AJ_0(z)+FJ_1(z),\\
 P_A(z)&=\nu_{t,A}(z)
       \left[\frac{t_A+FL_h(z)}{D_A(z)}\right]^2.
 \end{aligned}
 \label{eq:accepted-forced-awf}
\end{equation}
Thus no slip, the original forcing cutoff at $h$, and the explicit
pressure-gradient response retain their analytical form. Their
numerical values change with the new viscosity coefficients.
The vector form follows \cref{eq:3d}.

At the counterfactual zero-forcing reference, choose
\begin{equation}
 U_{H,0}=t_0J_0(H),\qquad
 U_{w,0}(z)=t_0J_0(z),\qquad
 p_{A,0}(z)=t_0^2\frac{\nu_{t,A}(z)}{D_A(z)^2}.
 \label{eq:accepted-reference-production}
\end{equation}
The resulting integrated reference production is
\begin{equation}
 \boxed{\mathcal P_0:=\int_0^H p_{A,0}(z)\,\dd z
 =t_0^2[J_0(H)-\nu J_2(H)]=HP_0.}
 \label{eq:accepted-production-budget}
\end{equation}
$P_0$ remains the height mean used in the main note; $p_{A,0}$ is
pointwise reference production. Neither is the production of a
general current terrain flow.

For $z\geq H$, preserve the standard outer CLS branch
\begin{equation}
 \boxed{\begin{aligned}
 U_{o,0}(z)&=t_0J_0(H)+\frac{\ust}{\kc}\ln\frac zH,\\
 k_{o,0}(z)&=K,\qquad
 \eps_{o,0}(z)=\frac{\Aref}{z},\qquad
 \nu_{t,o,0}(z)=\kc\ust z.
 \end{aligned}}
 \label{eq:accepted-outer-reference}
\end{equation}
The outer velocity offset has been used to satisfy the AWF anchor.
Its aerodynamic intercept is therefore
\begin{equation}
 z_{0,c}=H\exp[-\kc\ust J_0(H)].
 \label{eq:accepted-aerodynamic-intercept}
\end{equation}
Retaining the same sand-grain $h$ while changing the AWF coefficients
generally changes this intercept. A separately prescribed $z_{0,c}$
requires an additional momentum/intercept calibration.

\subsection{Budget-determined dissipation and admissibility}
\label{app:accepted-awf-cutoff}

Use the continuous positive dissipation distribution
\begin{equation}
 \boxed{E_w(z)=\ECw(z)=\frac{\Aref}{\max(z,z_c)},
 \qquad C_k^w(z)=0.}
 \label{eq:accepted-dissipation}
\end{equation}
The required tail begins no later than $H$, so $0<z_c\leq H$.
Its integral is
\begin{equation}
 \int_0^H\ECw(z)\,\dd z
 =\Aref\left(1+\ln\frac H{z_c}\right).
 \label{eq:accepted-dissipation-integral}
\end{equation}
The standard outer reference has zero $k$ gradient. Together with
zero modeled ground $k$ flux, conservative matching therefore
requires
\begin{equation}
 \boxed{\int_0^H\ECw\,\dd z=\mathcal P_0.}
 \label{eq:accepted-budget-condition}
\end{equation}
Combining the preceding identities gives
\begin{equation}
 \boxed{z_c=H\exp(1-\Lambda),\qquad
 \Lambda=\frac{\mathcal P_0}{\Aref}
 =\frac{t_0^2[J_0(H)-\nu J_2(H)]}{\Aref}.}
 \label{eq:accepted-cutoff}
\end{equation}
The left side of \cref{eq:accepted-dissipation-integral} decreases
strictly with $z_c$, tends to infinity as $z_c\downarrow0$, and equals
$\Aref$ at $z_c=H$. Consequently,
\begin{equation}
 \boxed{0<z_c\leq H\quad\Longleftrightarrow\quad\Lambda\geq1,}
 \label{eq:accepted-cutoff-gate}
\end{equation}
and the admissible cutoff is unique. This is the $z_m=H$
specialization of \cref{eq:hp-gate}. If the gate fails, this
monotone cutoff family cannot meet both the budget and outer trace.
Clipping $z_c$ to $H$ would not preserve the budget.

The original molecular cutoff
$z_{\eps,A}=2c_\ell^w\nu/\sqrt K$ is replaced by
\cref{eq:accepted-cutoff}. The plateau is now $\Aref/z_c$.
It cannot simultaneously remain $2\nu K/z_c^2$: imposing both
plateau formulas would force the original molecular cutoff again.
Thus the cutoff change is an explicit change of the inner
dissipation model.

\subsection{Derivation of the nonconstant inner reference}
\label{app:accepted-awf-inner-k}

Retain the inner transport equation with its physical production
and prescribed dissipation,
\begin{equation}
 \partial_t k_w+\bm U_w\cdot\nabla k_w
 =\nabla\cdot(D_{k,w}\nabla k_w)+P_A-\ECw.
 \label{eq:accepted-k-transport}
\end{equation}
For the steady, parallel flat reference this reduces to
\begin{equation}
 (D_{k,w}k_{w,0}')'=\ECw-p_{A,0},\qquad
 D_{k,w}(0)k_{w,0}'(0)=0,\qquad k_{w,0}(H)=K.
 \label{eq:accepted-k-bvp}
\end{equation}
Define the cumulative source defect
\begin{equation}
 Q_0(z)=\int_0^z[\ECw(s)-p_{A,0}(s)]\,\dd s.
 \label{eq:accepted-Q}
\end{equation}
Two integrations give the unique solution for the prescribed
positive diffusivity,
\begin{equation}
 \boxed{D_{k,w}(z)k_{w,0}'(z)=Q_0(z),\qquad
 k_{w,0}(z)=K-\int_z^H\frac{Q_0(s)}{D_{k,w}(s)}\,\dd s.}
 \label{eq:accepted-k-profile}
\end{equation}
The quantity $Q_0$ is the gradient flux inside the diffusion
operator; the physical upward diffusive flux is $-Q_0$.
The cutoff budget gives
\begin{equation}
 Q_0(0)=Q_0(H)=0,\qquad k_{w,0}'(H)=0.
 \label{eq:accepted-endpoint-fluxes}
\end{equation}
Conversely, these two zero endpoint fluxes imply
\cref{eq:accepted-budget-condition} by integration of the $k$
equation. Allowing a variable $k_w$ supplies local diffusion; it
does not remove the integrated compatibility requirement.

No additional value of $k_w(0)$ is prescribed. It follows from
\cref{eq:accepted-k-profile}. This is a modeled rough-wall layer,
not a viscous-sublayer resolution with an independently imposed
$k(0)=0$. The formula defines a stationary reference and an
initialization; during a general flow, $k_w$ is evolved by
\cref{eq:accepted-k-transport}, rather than reset to this profile.

For every sufficiently regular test function $\phi$, integration
by parts yields
\begin{equation}
 \boxed{-\int_0^H D_{k,w}k_{w,0}'\phi'\,\dd z
 +\int_0^H(p_{A,0}-\ECw)\phi\,\dd z=0.}
 \label{eq:accepted-weak-equilibrium}
\end{equation}
Thus the nonconstant reference balances every continuum test
moment. The higher-mode obstruction in \cref{sec:hp-source}
assumed constant inner $k$, for which the diffusion term vanished.
Finite-dimensional projection and numerical differentiation can
still leave a discrete residual; continuum balance alone does not
establish exact spectral/$hp$ preservation.

\subsection{Closed formulas and positivity on the rough branch}
\label{app:accepted-awf-rough}

Suppose $z_v\leq0$ and define
\begin{equation}
 d=\nu-az_v\geq\nu,\qquad D_A(z)=d+az.
 \label{eq:accepted-rough-D}
\end{equation}
Then
\begin{equation}
 \begin{aligned}
 J_0(z)&=\frac1a\ln\frac{d+az}{d},&
 J_2(z)&=\frac1a\left(\frac1d-\frac1{d+az}\right),\\
 p_{A,0}(z)&=t_0^2\frac{d+az-\nu}{(d+az)^2}.&&
 \end{aligned}
 \label{eq:accepted-rough-integrals}
\end{equation}
Using $\Aref=t_0^2/a$, the cumulative production is
\begin{equation}
 \int_0^z p_{A,0}(s)\,\dd s
 =\Aref\left[
 \ln\frac{d+az}{d}
 -\nu\left(\frac1d-\frac1{d+az}\right)\right].
 \label{eq:accepted-rough-production-integral}
\end{equation}
In particular, with $\zeta_H=aH/d$ and $\beta=\nu/d$,
\begin{equation}
 \Lambda=\ln(1+\zeta_H)
          -\frac{\beta\zeta_H}{1+\zeta_H}.
 \label{eq:accepted-rough-Lambda}
\end{equation}
The cumulative dissipation is $\Aref R_c(z)$, where
\begin{equation}
 R_c(z)=\begin{cases}
 z/z_c,&0\leq z\leq z_c,\\
 1+\ln(z/z_c),&z_c\leq z\leq H.
 \end{cases}
 \label{eq:accepted-Rc}
\end{equation}
Hence an elementary expression for the flux is
\begin{equation}
 Q_0(z)=\Aref\left[
 R_c(z)-\ln\frac{d+az}{d}
 +\nu\left(\frac1d-\frac1{d+az}\right)\right].
 \label{eq:accepted-rough-Q}
\end{equation}
Together with \cref{eq:accepted-k-profile}, this provides an exact
analytical representation of the inner reference.

\paragraph{A closed primitive for the kinetic energy.}
For the selected $\sigma_k=1$, set
$\zeta=az/d$, $\zeta_c=az_c/d$, and define
\begin{align}
 \mathcal F_< (\zeta)
 &=\frac{\zeta-\ln(1+\zeta)}{\zeta_c}
   -\frac12\ln^2(1+\zeta)
   +\beta\left[\ln(1+\zeta)+\frac1{1+\zeta}-1\right],
 \label{eq:accepted-primitive-lower}\\
 \mathcal G(\zeta)
 &=\left[1+\ln\frac{\zeta}{\zeta_c}+\beta\right]\ln(1+\zeta)
   +\operatorname{Li}_2(-\zeta)
   -\frac12\ln^2(1+\zeta)+\frac{\beta}{1+\zeta},
 \quad \zeta>0.
 \label{eq:accepted-primitive-upper}
\end{align}
The real dilogarithm is defined here by
$\operatorname{Li}_2(-\zeta)=-\int_0^\zeta\ln(1+s)\,\dd s/s$.
Join the primitives continuously:
\begin{equation}
 \mathcal F(\zeta)=
 \begin{cases}
 \mathcal F_< (\zeta),&0\leq\zeta\leq\zeta_c,\\
 \mathcal F_< (\zeta_c)+\mathcal G(\zeta)-\mathcal G(\zeta_c),
 &\zeta_c\leq\zeta\leq\zeta_H.
 \end{cases}
 \label{eq:accepted-primitive-joined}
\end{equation}
Direct differentiation of these expressions gives
$\mathcal F'(\zeta)=Q_0(d\zeta/a)/[\Aref(1+\zeta)]$.
Therefore the closed profile is
\begin{equation}
 \boxed{k_{w,0}(z)=K+\frac{\Aref}{a}
 \left[\mathcal F\!\left(\frac{az}{d}\right)
       -\mathcal F(\zeta_H)\right].}
 \label{eq:accepted-k-closed}
\end{equation}
The integral form \cref{eq:accepted-k-profile} remains applicable
to other viscosity segments and positive diffusivities.

\paragraph{The sign of the local imbalance.}
For every $z\geq z_c>0$ on this rough branch,
\begin{equation}
 \boxed{\ECw(z)-p_{A,0}(z)
 =\frac{\Aref[(d+\nu)az+d^2]}{z(d+az)^2}>0.}
 \label{eq:accepted-tail-defect}
\end{equation}
This strict inequality concerns the $1/z$ tail. It does not imply
positive imbalance throughout the plateau. In fact, the budget
requires a compensating region of negative imbalance.

\begin{proposition}[A positive rough-wall reference]
Assume $z_v\leq0$, $d\geq2\nu$, and $\Lambda\geq1$.
Then $\ECw-p_{A,0}$ has exactly one zero in $(0,z_c)$,
changing from negative to positive. For every positive $D_{k,w}$,
\begin{equation}
 Q_0(z)<0\quad(0<z<H),\qquad
 k_{w,0}'(z)<0\quad(0<z<H),\qquad
 k_{w,0}(z)>K>0\quad(0\leq z<H).
 \label{eq:accepted-positivity}
\end{equation}
\end{proposition}
\begin{proof}
The reference production satisfies
\[
 p_{A,0}'(z)=
 \frac{a t_0^2[2\nu-(d+az)]}{(d+az)^3}\leq0,
\]
and is strictly decreasing for $z>0$. On the dissipation plateau,
$\Aref/z_c-p_{A,0}$ is therefore strictly increasing and is
positive at $z_c$ by \cref{eq:accepted-tail-defect}.
Its initial value must be negative, since otherwise the total
integral could not vanish. It crosses zero exactly once, after
which the defect remains positive up to $H$. Its cumulative
integral first decreases and then increases to $Q_0(H)=0$, so
$Q_0<0$ in the interior. The profile and monotonicity statements
follow from \cref{eq:accepted-k-profile}.
\end{proof}
The condition $d\geq2\nu$ is sufficient, not necessary.
For other intercepts, including a smooth-wall viscosity segment,
the budget construction remains defined when its gate holds, but
positivity must be assessed from the resulting profile.

The rough-branch velocity anchor also gives the useful identities
\begin{equation}
 \boxed{z_{0,c}=\frac{dH}{aH+d},\qquad
 z_c=\ee z_{0,c}
 \exp\!\left[\nu\left(\frac1d-\frac1{aH+d}\right)\right].}
 \label{eq:accepted-roughness-cutoff}
\end{equation}
Even with the calibrated slope, $z_{0,c}$ retains a finite-height
dependence. For $aH\gg d$, it tends to $d/a$.

\subsection{Conservative matching and normal pressure balance}
\label{app:accepted-awf-matching}

The inner momentum reference and the outer profile in
\cref{eq:accepted-outer-reference} satisfy
\begin{equation}
 \begin{gathered}
 U_{w,0}(0)=0,\qquad U_{w,0}(H)=U_{o,0}(H),\\
 D_A(H)U_{w,0}'(H)
 =\nu_{t,o,0}(H)U_{o,0}'(H)=t_0.
 \end{gathered}
 \label{eq:accepted-momentum-match}
\end{equation}
The calibrated coefficients match the viscosity slope, not the
finite-height viscosities: on the rough branch,
$D_A(H)=aH+d$ while $\nu_{t,o,0}(H)=aH$.
The velocity derivatives therefore need not coincide.

The turbulence matching conditions are
\begin{equation}
 \boxed{\begin{aligned}
 k_{w,0}(H)&=k_{o,0}(H)=K,\\
 D_{k,w}(H)k_{w,0}'(H)
 &=D_{k,o}(H)k_{o,0}'(H)=0,\\
 \ECw(H)&=\eps_{o,0}(H)=\frac{\Aref}{H}.
 \end{aligned}}
 \label{eq:accepted-turbulence-match}
\end{equation}
Here $D_{k,o}=\nu_t^{CLS}/\sigma_k$ on the declared outer branch.
Inner dissipation is dependent, so the outer dissipation equation
receives its single trace. No additional dissipation-flux condition
is imposed. In the limiting case $z_c=H$, even the one-sided
dissipation derivatives differ, while the required trace still
matches.

With the full-strain inner stress completion of
\cref{eq:curve-inner-stress}, the parallel inner normal stress is
$R_{zz,w}=2k_w/3$. Its normal momentum balance therefore requires
\begin{equation}
 \frac{\dd}{\dd z}\left(
 \frac{p_w}{\rho}+gz+\frac23 k_{w,0}\right)=0.
 \label{eq:accepted-inner-pressure}
\end{equation}
Mechanical pressure inside this layer is generally not hydrostatic.
Since $R_{zz,o}=0.52K$ at the outer reference, continuity of total
normal traction gives
\begin{equation}
 \boxed{\frac{p_w(H)-p_o(H)}{\rho}
 =\left(0.52-\frac23\right)K=-\frac{11}{75}K.}
 \label{eq:accepted-pressure-jump}
\end{equation}
The interface offset is unchanged from the full-strain EP-AWF
completion because the matching value is still $K$; the interior
pressure distribution changes with $k_{w,0}(z)$.
When a reduced pressure is stored, recover the mechanical pressure
consistently before evaluating the AWF pressure forcing.

The flat reference also retains the integrated mechanical-energy
identity
\begin{equation}
 t_0 U_{H,0}
 =\nu\int_0^H[U_{w,0}'(z)]^2\,\dd z
  +\int_0^H p_{A,0}(z)\,\dd z
 =\nu t_0^2J_2(H)+\int_0^H\ECw(z)\,\dd z.
 \label{eq:accepted-energy}
\end{equation}
The nonzero internal $k$ flux redistributes energy within the layer;
it contributes no net flux at the two reference endpoints.

\subsection{Distinction from EP-AWF with unit equilibrium weight}
\label{app:accepted-awf-versus-ep}

To isolate the transport difference, compare the two constructions
at identical $K,h,H,\nu$, viscosity coefficients, and budget cutoff.
At $\chi=1$, the EP-AWF has
\begin{equation}
 E_w^{EP}=\ECw,\qquad
 C_{k,EP}^w=\ECw-p_{A,0},\qquad
 S_{k,w}^{EP}=P_A-p_{A,0}.
 \label{eq:accepted-ep-comparison}
\end{equation}
At the flat reference, $P_A=p_{A,0}$ and this source vanishes
pointwise, permitting constant inner $k=K$. The accepted alternative
instead uses
\begin{equation}
 E_w=\ECw,\qquad C_k^w=0,\qquad
 S_{k,w}=P_A-\ECw.
 \label{eq:accepted-source-comparison}
\end{equation}
Its reference source is generally nonzero and is balanced by
$\partial_z(D_{k,w}k_{w,0}')$.

Both sources have the same layer integral,
\begin{equation}
 \int_0^H S_{k,w}\,\dd z
 =\int_0^H S_{k,w}^{EP}\,\dd z
 =\int_0^H(P_A-p_{A,0})\,\dd z,
 \label{eq:accepted-equal-integrated-sources}
\end{equation}
but their local balances and inner kinetic-energy fields differ.
On the reference, EP-AWF carries the physical redistribution flux
$-Q_0$ through its prescribed correction; the accepted model carries
the same flux through ordinary diffusion of its nonconstant $k$.
This reference identity does not make their evolving dynamics
equivalent. In addition, comparison with an EP-AWF that retains the
original coefficients must account separately for the changed
momentum reconstruction.

Setting the old switch to $\chi=0$ does not select the present
variant: it restores the original-cutoff dissipation $\EAw$ through
\cref{eq:blend}. The accepted model retains $\ECw$ while removing
$C_k^w$. Consequently the spatial fade to the original AWF in
\cref{sec:transition} is not part of this variant.

\subsection{A specified extension to curved terrain}
\label{app:accepted-awf-curved}

The accepted scalar policy can be used with the conservative curved
patch of \cref{sec:curved}. To make that extension definite, adopt
the note's metric-weighted reference budget, together with
$r_K=r_m=H$, $K=k|_{\Gamma_H}$, the calibrated viscosity, and
$C_k^w=0$. This is an explicit extension of the flat model; the
flat matching theorem does not by itself select a unique curved
calibration or establish an equilibrium over a hill.

Retain the plus curvature convention
\begin{equation}
 \bm x(\bm\xi,r)=\bm X(\bm\xi)+r\bm n(\bm\xi),\qquad
 \mathsf B=\nabla_\Gamma\bm n,\qquad
 \mathsf A=I_T+r\mathsf B,\qquad J=\det_T\mathsf A.
 \label{eq:accepted-curved-map}
\end{equation}
Here $r$ is physical normal distance and $\dd V=J\,\dd S_w\,\dd r$.
All principal stretches $1+r\kappa_i$ must stay positive, and
the normal map must be globally one-to-one. Positive local $J$
alone does not guarantee unique normal projection.

\paragraph{Reference cutoff in the physical volume measure.}
Define
\begin{equation}
 \boxed{\begin{aligned}
 E_w(r)&=\ECw(r)=\frac{\Aref(K)}{\max(r,r_c)},\\
 \int_0^H J(r)\ECw(r)\,\dd r
 &=\mathcal P_J:=\int_0^H J(r)p_{A,0}(r;K)\,\dd r,\\
 p_{A,0}(r;K)&=t_0(K)^2
              \frac{\nu_{t,A}(r;K)}{D_A(r;K)^2}.
 \end{aligned}}
 \label{eq:accepted-curved-budget}
\end{equation}
The reference kernel remains counterfactual flat, parallel,
zero-forcing production at the current $K$. It is not replaced
by actual terrain production.
For $\mathcal B(c)=\int_0^H J(r)\Aref/\max(r,c)\,\dd r$,
\begin{equation}
 \mathcal B'(c)=-\frac{\Aref}{c^2}\int_0^cJ(r)\,\dd r<0,
 \qquad \lim_{c\downarrow0}\mathcal B(c)=+\infty.
 \label{eq:accepted-curved-monotonicity}
\end{equation}
Therefore a unique $0<r_c\leq H$ exists exactly when
\begin{equation}
 \boxed{\mathcal P_J\geq
 \frac{\Aref}{H}\int_0^HJ(r)\,\dd r.}
 \label{eq:accepted-curved-gate}
\end{equation}
At a fixed footpoint on a three-dimensional surface,
$J=(1+\kappa_1r)(1+\kappa_2r)$. Direct integration gives
\begin{equation}
 \begin{aligned}
 \frac{\mathcal B(c)}{\Aref}
 &=1+\ln\frac Hc
 +(\kappa_1+\kappa_2)\left(H-\frac c2\right)\\
 &\quad+\kappa_1\kappa_2
       \left(\frac{H^2}{2}-\frac{c^2}{6}\right).
 \end{aligned}
 \label{eq:accepted-curved-budget-explicit}
\end{equation}
The curved cutoff thus requires one monotone scalar root solve.
The exponential formula is recovered when $J=1$.

\paragraph{Actual inner flow and transport.}
Write $\bm U=\bm u_t+w\bm n$ and
$\bm v=\mathsf A^{-1}\bm u_t$. The inner equation is
\begin{equation}
 \boxed{\begin{aligned}
 J\partial_t k
 &+\partial_r\{J(wk-D_{k,w}\partial_rk)\}\\
 &+\operatorname{div}_\Gamma\left\{
 J[\bm v k-D_{k,w}\mathsf A^{-1}\mathsf A^{-T}
                         \nabla_\Gamma k]\right\}
 =J(P_A-\ECw).
 \end{aligned}}
 \label{eq:accepted-curved-k}
\end{equation}
Use the full continuity and momentum equations
\cref{eq:curve-continuity,eq:curve-momentum}, the stated roughness
force in \cref{eq:curve-rough-force}, and the inner completion
\begin{equation}
 \begin{gathered}
 \R_A=\tfrac23k\I-\nu_{t,A}\Sm_w,\qquad
 \mathsf T_w=-\left(\frac{p_w}{\rho}+\tfrac23k\right)\I
              +D_A\Sm_w,\\
 P_A=-\R_A:\nabla\bm U
     =\tfrac12\nu_{t,A}\Sm_w:\Sm_w
 \quad\text{for }\nabla\cdot\bm U=0.
 \end{gathered}
 \label{eq:accepted-curved-stress}
\end{equation}
The isotropic stress uses local transported $k$, not sampled $K$.
For a weakly incompressible discrete field, the production
convention discussed after \cref{eq:curve-inner-stress} retains
the corresponding $-(2k/3)\nabla\cdot\bm U$ contribution when
required by the discrete stress-work balance.

The base-normal tangential traction contains curvature and normal
motion,
\begin{equation}
 \bm t_A=D_A\left[
 \partial_r\bm u_t+\mathsf A^{-T}\nabla_\Gamma w
 -\mathsf B\mathsf A^{-1}\bm u_t\right].
 \label{eq:accepted-curved-traction}
\end{equation}
Normal momentum must likewise retain curvature acceleration and
pressure balance. Merely weighting independent flat momentum
columns by $J$ does not supply these terms. General curved and
developing flow uses the coupled patch solve of \cref{sec:curved}.

The reference budget preserves the actual non-equilibrium source:
\begin{equation}
 \int_0^H J(P_A-\ECw)\,\dd r
 =\int_0^H J(P_A-p_{A,0})\,\dd r.
 \label{eq:accepted-curved-actual-budget}
\end{equation}
This quantity need not vanish; storage, tangential transport, and
interface flux balance it. Determining the cutoff from actual
$P_A$ instead of $p_{A,0}$ would impose a different energy budget.

An auxiliary frozen normal-column construction can still be written,
\begin{equation}
 \begin{aligned}
 Q_J(r)&=\int_0^rJ(s)[\ECw(s)-p_{A,0}(s)]\,\dd s,\\
 k_{\rm col}(r)&=K-\int_r^H
            \frac{Q_J(s)}{J(s)D_{k,w}(s)}\,\dd s.
 \end{aligned}
 \label{eq:accepted-curved-column}
\end{equation}
It balances only the reference source and normal diffusion.
It is not an exact terrain-flow equilibrium, since actual
production, advection, and tangential diffusion generally differ.
Its zero base-normal endpoint flux also does not establish zero
flux through an arbitrary sloping matching surface.

\paragraph{Interface conditions.}
Retain the common velocity, total traction, and mass flux in
\cref{eq:curve-momentum-interface,eq:curve-massflux}, together with
the common $k$ trace and advective--diffusive flux in
\cref{eq:curve-k-interface}. For varying $H(\bm\xi)$, these use
\begin{equation}
 \bm N_H=\frac{\bm n-\mathsf A_H^{-T}\nabla_\Gamma H}
 {\sqrt{1+|\mathsf A_H^{-T}\nabla_\Gamma H|^2}}.
 \label{eq:accepted-curved-interface-normal}
\end{equation}
Surface derivatives in the mapped equations are taken at fixed
$r$ before restriction to $H$; in particular,
$\nabla_\Gamma K=(\nabla_\Gamma k)|_H
+(\partial_rk)|_H\nabla_\Gamma H$.
The outer dissipation equation receives only
\begin{equation}
 \eps_{o,H}=\ECw(H)=\frac{\Aref(K)}{H}.
 \label{eq:accepted-curved-epsilon-trace}
\end{equation}
The full normal-traction condition replaces the flat pressure
offset over terrain. There is no independently solved inner
dissipation equation or additional prescribed dissipation flux.

\subsection{Illustrative reference and implementation implications}
\label{app:accepted-awf-example}

For a flat reference with
\begin{equation}
 \ust=0.4\units{m\,s^{-1}},\quad
 \nu=1.5\times10^{-5}\units{m^2\,s^{-1}},\quad
 h=0.1\units{m},\quad H=2\units{m},\quad z_K=z_m=H,
 \label{eq:accepted-example-inputs}
\end{equation}
the unchanged roughness correlation and calibrated coefficients give
the values in \cref{tab:accepted-reference}. They evaluate the
analytical construction; they are not a solver-validation result.

\begin{table}[htbp]
\centering\small
\begin{tabular}{@{}lll@{}}
\toprule
Quantity & Value & Units\\
\midrule
$K$ & $0.5105080131$ & $\mathrm{m^2\,s^{-2}}$\\
$z_v$ & $-0.0052721312$ & $\mathrm m$\\
$a$ & $0.1762215621$ & $\mathrm{m\,s^{-1}}$\\
$d$ & $0.0009440632$ & $\mathrm{m^2\,s^{-1}}$\\
$\mathcal P_0$ & $0.8584509755$ & $\mathrm{m^3\,s^{-3}}$\\
$z_c$ & $0.0147555946$ & $\mathrm m$\\
$U_{H,0}$ & $5.3797062304$ & $\mathrm{m\,s^{-1}}$\\
$z_{0,c}$ & $0.0053429396$ & $\mathrm m$\\
$k_{w,0}(0)$ & $1.0229104395$ & $\mathrm{m^2\,s^{-2}}$\\
$k_{w,0}(0.1\units m)$ & $0.5447982810$ & $\mathrm{m^2\,s^{-2}}$\\
$k_{w,0}(1\units m)$ & $0.5111909311$ & $\mathrm{m^2\,s^{-2}}$\\
\bottomrule
\end{tabular}
\caption{Accepted flat-wall reference with endpoint sampling.
The outer kinetic energy remains $K$, while ordinary diffusion
supports the larger inner values.}
\label{tab:accepted-reference}
\end{table}

Here $d>2\nu$, so the positivity proposition applies. In particular,
$\ECw(0)\simeq9.84519\units{m^2\,s^{-3}}$ is smaller than
$p_{A,0}(0)\simeq26.68598\units{m^2\,s^{-3}}$.
The defect changes sign near $z=0.00931273\units m$ and remains
positive thereafter, while its full-layer integral is zero.
These signs refer to the accepted budget cutoff; retaining the
original molecular cutoff would give a different imbalance.

Implementation should evaluate $K$ from the common matching trace,
construct $D_A$ with the calibrated coefficients, and determine
the cutoff from reference production at that same $K$.
Actual production is used in \cref{eq:accepted-k-transport} or
\cref{eq:accepted-curved-k}. Initialize a flat reference with
\cref{eq:accepted-k-profile}, retain zero modeled ground $k$ flux,
and couple the matching value and conservative flux to outer CLS.
Do not reintroduce the EP redistribution while using this profile.
Quadrature should resolve each active viscosity, forcing, and
dissipation breakpoint, and the full assembled reference residual
must be checked with the actual interface and pressure operators.

The derivation establishes continuum compatibility with the
standard outer CLS equilibrium on flat, undisturbed terrain and
specifies a conservative curved extension. It does not require
the outer flow over a hill to remain logarithmic. The positive
flat reference is also not a positivity or stability theorem for
the evolving equations: a sink controlled by sampled $K$ can
remain positive as local $k_w$ decreases, and the constitutive
limitations of the unchanged outer CLS model identified in
\cref{sec:stability} remain. Any active robustness modification
must be assessed against this reference and declared separately.

% The main note already calls \appendix. Do not call it again here.
% Uses the main note's macros, theorem environments, labels and bibliography.

\section{AWF with CLS-compatible coefficients, a budget-determined cutoff,
and a nonconstant inner kinetic-energy profile}
\label{app:accepted-awf}

This appendix defines the accepted alternative to the constant-inner-$k$
EP-AWF construction. The objective is to preserve the standard neutral
CLS equilibrium in the outer region while allowing the kinetic energy
to vary inside the modeled wall layer. The AWF viscosity coefficients
are calibrated to that outer equilibrium, dissipation retains the
budget-determined cutoff, and the prescribed redistribution $C_k^w$
is removed. Ordinary diffusion of the transported inner $k$ then
balances the local difference between production and dissipation.
The outer CLS model, including its Yap correction, is unchanged.

The original CLS and rough-wall AWF supply the constitutive and
wall-model starting points \cite{CLS1996,SCI2006}. The coefficient
calibration, budget cutoff, retained inner transport, and curved
completion specified here are modifications of that combination.
When this variant is selected, the definitions below replace the
original wall coefficients, interior sampling choice, and
$\chi$-dependent source policy used elsewhere in the note.

\subsection{Scope, endpoint sampling, and coefficient calibration}
\label{app:accepted-awf-coefficients}

Use the high-Re outer equation convention of
\cref{sec:derivation}: molecular momentum and scalar diffusion are
omitted in the outer analytical reference, while $\nu>0$ is retained
in the AWF. This is a statement about the selected equations, not an
assertion that the same profiles solve a finite-Re outer system.
The transported inner turbulence variable is $k$; inner
dissipation is a dependent total $k$ sink. No inner CLS dissipation
transport equation is introduced.

The accepted sampling and tail-matching heights are
\begin{equation}
 \boxed{z_K=z_m=H,\qquad
 K=k_w(H)=k_o(H)>0.}
 \label{eq:accepted-sampling}
\end{equation}
The subscripts $w,o$ distinguish the modeled wall layer and outer
region. At each positive current $K$, define the counterfactual
reference scales
\begin{equation}
 \ust(K)=\sqrt{\frac{K}{\qs}},\qquad
 t_0(K)=\frac{K}{\qs},\qquad
 \Aref(K)=\frac{\ust(K)^3}{\kc}.
 \label{eq:accepted-reference-scales}
\end{equation}
These are reference scales. In a developing, forced, or separated
flow, $t_0(K)$ is not imposed as the actual wall traction.
Arguments $K$ are suppressed below when all quantities use the same
sample. A fixed inlet reference is obtained by setting $K=K_0$.

The outer constants are the roots derived in
\cref{eq:qroot,eq:kappaeq},
\begin{equation}
 \begin{gathered}
 C(\qs)\qs^2=1,\qquad
 \frac{\qs\kc^2}{1.3}
 +H_Y\!\left(\frac{\kc\qs^{3/2}}{2.5}\right)=0.48,\\
 \qs=3.190675081624288,\qquad
 \kc=0.440553905186659,\qquad
 \Cs=\qs^{-2}=0.0982278953282113.
 \end{gathered}
 \label{eq:accepted-cls-constants}
\end{equation}
Here $C$ and $H_Y$ retain the definitions of the original high-Re CLS
closure and its Yap correction. Matching the dissipation coefficient
and the asymptotic auxiliary-viscosity slope gives
\begin{equation}
 \boxed{\begin{aligned}
 c_\ell^w&=\kc\qs^{3/2}=2.510862475219797,\\
 c_\mu^w&=\Cs=\qs^{-2}=0.0982278953282113,\\
 \alpha&=c_\mu^wc_\ell^w
       =\frac{\kc}{\sqrt{\qs}}=0.246636736399424.
 \end{aligned}}
 \label{eq:accepted-coefficients}
\end{equation}
Indeed, with $a:=\alpha\sqrt K$,
\begin{equation}
 \boxed{a=\kc\ust,\qquad
 \frac{K^{3/2}}{c_\ell^w}=\Aref
 =\frac{t_0^2}{a}.}
 \label{eq:accepted-calibration-identities}
\end{equation}
The choice of $c_\ell^w$ matches the outer dissipation tail exactly.
The additional choice of $c_\mu^w$ matches the viscosity slope;
conservative interface matching by itself does not require equality
of these slopes.

Retain the original intercept correlation $b(K,h,\nu)$ in
\cref{eq:brough}, including $b=10.7$ for $h=0$. Define
\begin{equation}
 z_v=\frac{\nu b}{\sqrt K},\qquad
 \nu_{t,A}(z)=a\max(z-z_v,0),\qquad
 D_A=\nu+\nu_{t,A},\qquad
 D_{k,w}=\nu+\frac{\nu_{t,A}}{\sigma_k}.
 \label{eq:accepted-viscosities}
\end{equation}
The CLS choice $\sigma_k=1$ gives $D_{k,w}=D_A$.
These auxiliary coefficients depend on sampled $K$, while the
transported field $k_w(z)$ can vary. Substituting $k_w(z)$ for $K$
in the viscosity or dissipation law would define a different model.

\subsection{Momentum, production, and the outer velocity anchor}
\label{app:accepted-awf-momentum}

The analytical form of the AWF momentum reconstruction is retained,
with all integrals evaluated using the calibrated $D_A$. Put
$L_h(z)=\max(z-h,0)$ and
\begin{equation}
 J_0(z)=\int_0^z\frac{\dd s}{D_A(s)},\qquad
 J_1(z)=\int_0^z\frac{L_h(s)}{D_A(s)}\,\dd s,\qquad
 J_2(z)=\int_0^z\frac{\dd s}{D_A(s)^2}.
 \label{eq:accepted-J}
\end{equation}
For scalar tangential forcing $F$, the flat quasi-steady reduction is
\begin{equation}
 \begin{aligned}
 t_A&=\frac{U_H-FJ_1(H)}{J_0(H)},\qquad
 U_A(z)=t_AJ_0(z)+FJ_1(z),\\
 P_A(z)&=\nu_{t,A}(z)
       \left[\frac{t_A+FL_h(z)}{D_A(z)}\right]^2.
 \end{aligned}
 \label{eq:accepted-forced-awf}
\end{equation}
Thus no slip, the original forcing cutoff at $h$, and the explicit
pressure-gradient response retain their analytical form. Their
numerical values change with the new viscosity coefficients.
The vector form follows \cref{eq:3d}.

At the counterfactual zero-forcing reference, choose
\begin{equation}
 U_{H,0}=t_0J_0(H),\qquad
 U_{w,0}(z)=t_0J_0(z),\qquad
 p_{A,0}(z)=t_0^2\frac{\nu_{t,A}(z)}{D_A(z)^2}.
 \label{eq:accepted-reference-production}
\end{equation}
The resulting integrated reference production is
\begin{equation}
 \boxed{\mathcal P_0:=\int_0^H p_{A,0}(z)\,\dd z
 =t_0^2[J_0(H)-\nu J_2(H)]=HP_0.}
 \label{eq:accepted-production-budget}
\end{equation}
$P_0$ remains the height mean used in the main note; $p_{A,0}$ is
pointwise reference production. Neither is the production of a
general current terrain flow.

For $z\geq H$, preserve the standard outer CLS branch
\begin{equation}
 \boxed{\begin{aligned}
 U_{o,0}(z)&=t_0J_0(H)+\frac{\ust}{\kc}\ln\frac zH,\\
 k_{o,0}(z)&=K,\qquad
 \eps_{o,0}(z)=\frac{\Aref}{z},\qquad
 \nu_{t,o,0}(z)=\kc\ust z.
 \end{aligned}}
 \label{eq:accepted-outer-reference}
\end{equation}
The outer velocity offset has been used to satisfy the AWF anchor.
Its aerodynamic intercept is therefore
\begin{equation}
 z_{0,c}=H\exp[-\kc\ust J_0(H)].
 \label{eq:accepted-aerodynamic-intercept}
\end{equation}
Retaining the same sand-grain $h$ while changing the AWF coefficients
generally changes this intercept. A separately prescribed $z_{0,c}$
requires an additional momentum/intercept calibration.

\subsection{Budget-determined dissipation and admissibility}
\label{app:accepted-awf-cutoff}

Use the continuous positive dissipation distribution
\begin{equation}
 \boxed{E_w(z)=\ECw(z)=\frac{\Aref}{\max(z,z_c)},
 \qquad C_k^w(z)=0.}
 \label{eq:accepted-dissipation}
\end{equation}
The required tail begins no later than $H$, so $0<z_c\leq H$.
Its integral is
\begin{equation}
 \int_0^H\ECw(z)\,\dd z
 =\Aref\left(1+\ln\frac H{z_c}\right).
 \label{eq:accepted-dissipation-integral}
\end{equation}
The standard outer reference has zero $k$ gradient. Together with
zero modeled ground $k$ flux, conservative matching therefore
requires
\begin{equation}
 \boxed{\int_0^H\ECw\,\dd z=\mathcal P_0.}
 \label{eq:accepted-budget-condition}
\end{equation}
Combining the preceding identities gives
\begin{equation}
 \boxed{z_c=H\exp(1-\Lambda),\qquad
 \Lambda=\frac{\mathcal P_0}{\Aref}
 =\frac{t_0^2[J_0(H)-\nu J_2(H)]}{\Aref}.}
 \label{eq:accepted-cutoff}
\end{equation}
The left side of \cref{eq:accepted-dissipation-integral} decreases
strictly with $z_c$, tends to infinity as $z_c\downarrow0$, and equals
$\Aref$ at $z_c=H$. Consequently,
\begin{equation}
 \boxed{0<z_c\leq H\quad\Longleftrightarrow\quad\Lambda\geq1,}
 \label{eq:accepted-cutoff-gate}
\end{equation}
and the admissible cutoff is unique. This is the $z_m=H$
specialization of \cref{eq:hp-gate}. If the gate fails, this
monotone cutoff family cannot meet both the budget and outer trace.
Clipping $z_c$ to $H$ would not preserve the budget.

The original molecular cutoff
$z_{\eps,A}=2c_\ell^w\nu/\sqrt K$ is replaced by
\cref{eq:accepted-cutoff}. The plateau is now $\Aref/z_c$.
It cannot simultaneously remain $2\nu K/z_c^2$: imposing both
plateau formulas would force the original molecular cutoff again.
Thus the cutoff change is an explicit change of the inner
dissipation model.

\subsection{Derivation of the nonconstant inner reference}
\label{app:accepted-awf-inner-k}

Retain the inner transport equation with its physical production
and prescribed dissipation,
\begin{equation}
 \partial_t k_w+\bm U_w\cdot\nabla k_w
 =\nabla\cdot(D_{k,w}\nabla k_w)+P_A-\ECw.
 \label{eq:accepted-k-transport}
\end{equation}
For the steady, parallel flat reference this reduces to
\begin{equation}
 (D_{k,w}k_{w,0}')'=\ECw-p_{A,0},\qquad
 D_{k,w}(0)k_{w,0}'(0)=0,\qquad k_{w,0}(H)=K.
 \label{eq:accepted-k-bvp}
\end{equation}
Define the cumulative source defect
\begin{equation}
 Q_0(z)=\int_0^z[\ECw(s)-p_{A,0}(s)]\,\dd s.
 \label{eq:accepted-Q}
\end{equation}
Two integrations give the unique solution for the prescribed
positive diffusivity,
\begin{equation}
 \boxed{D_{k,w}(z)k_{w,0}'(z)=Q_0(z),\qquad
 k_{w,0}(z)=K-\int_z^H\frac{Q_0(s)}{D_{k,w}(s)}\,\dd s.}
 \label{eq:accepted-k-profile}
\end{equation}
The quantity $Q_0$ is the gradient flux inside the diffusion
operator; the physical upward diffusive flux is $-Q_0$.
The cutoff budget gives
\begin{equation}
 Q_0(0)=Q_0(H)=0,\qquad k_{w,0}'(H)=0.
 \label{eq:accepted-endpoint-fluxes}
\end{equation}
Conversely, these two zero endpoint fluxes imply
\cref{eq:accepted-budget-condition} by integration of the $k$
equation. Allowing a variable $k_w$ supplies local diffusion; it
does not remove the integrated compatibility requirement.

No additional value of $k_w(0)$ is prescribed. It follows from
\cref{eq:accepted-k-profile}. This is a modeled rough-wall layer,
not a viscous-sublayer resolution with an independently imposed
$k(0)=0$. The formula defines a stationary reference and an
initialization; during a general flow, $k_w$ is evolved by
\cref{eq:accepted-k-transport}, rather than reset to this profile.

For every sufficiently regular test function $\phi$, integration
by parts yields
\begin{equation}
 \boxed{-\int_0^H D_{k,w}k_{w,0}'\phi'\,\dd z
 +\int_0^H(p_{A,0}-\ECw)\phi\,\dd z=0.}
 \label{eq:accepted-weak-equilibrium}
\end{equation}
Thus the nonconstant reference balances every continuum test
moment. The higher-mode obstruction in \cref{sec:hp-source}
assumed constant inner $k$, for which the diffusion term vanished.
Finite-dimensional projection and numerical differentiation can
still leave a discrete residual; continuum balance alone does not
establish exact spectral/$hp$ preservation.

\subsection{Closed formulas and positivity on the rough branch}
\label{app:accepted-awf-rough}

Suppose $z_v\leq0$ and define
\begin{equation}
 d=\nu-az_v\geq\nu,\qquad D_A(z)=d+az.
 \label{eq:accepted-rough-D}
\end{equation}
Then
\begin{equation}
 \begin{aligned}
 J_0(z)&=\frac1a\ln\frac{d+az}{d},&
 J_2(z)&=\frac1a\left(\frac1d-\frac1{d+az}\right),\\
 p_{A,0}(z)&=t_0^2\frac{d+az-\nu}{(d+az)^2}.&&
 \end{aligned}
 \label{eq:accepted-rough-integrals}
\end{equation}
Using $\Aref=t_0^2/a$, the cumulative production is
\begin{equation}
 \int_0^z p_{A,0}(s)\,\dd s
 =\Aref\left[
 \ln\frac{d+az}{d}
 -\nu\left(\frac1d-\frac1{d+az}\right)\right].
 \label{eq:accepted-rough-production-integral}
\end{equation}
In particular, with $\zeta_H=aH/d$ and $\beta=\nu/d$,
\begin{equation}
 \Lambda=\ln(1+\zeta_H)
          -\frac{\beta\zeta_H}{1+\zeta_H}.
 \label{eq:accepted-rough-Lambda}
\end{equation}
The cumulative dissipation is $\Aref R_c(z)$, where
\begin{equation}
 R_c(z)=\begin{cases}
 z/z_c,&0\leq z\leq z_c,\\
 1+\ln(z/z_c),&z_c\leq z\leq H.
 \end{cases}
 \label{eq:accepted-Rc}
\end{equation}
Hence an elementary expression for the flux is
\begin{equation}
 Q_0(z)=\Aref\left[
 R_c(z)-\ln\frac{d+az}{d}
 +\nu\left(\frac1d-\frac1{d+az}\right)\right].
 \label{eq:accepted-rough-Q}
\end{equation}
Together with \cref{eq:accepted-k-profile}, this provides an exact
analytical representation of the inner reference.

\paragraph{A closed primitive for the kinetic energy.}
For the selected $\sigma_k=1$, set
$\zeta=az/d$, $\zeta_c=az_c/d$, and define
\begin{align}
 \mathcal F_< (\zeta)
 &=\frac{\zeta-\ln(1+\zeta)}{\zeta_c}
   -\frac12\ln^2(1+\zeta)
   +\beta\left[\ln(1+\zeta)+\frac1{1+\zeta}-1\right],
 \label{eq:accepted-primitive-lower}\\
 \mathcal G(\zeta)
 &=\left[1+\ln\frac{\zeta}{\zeta_c}+\beta\right]\ln(1+\zeta)
   +\operatorname{Li}_2(-\zeta)
   -\frac12\ln^2(1+\zeta)+\frac{\beta}{1+\zeta},
 \quad \zeta>0.
 \label{eq:accepted-primitive-upper}
\end{align}
The real dilogarithm is defined here by
$\operatorname{Li}_2(-\zeta)=-\int_0^\zeta\ln(1+s)\,\dd s/s$.
Join the primitives continuously:
\begin{equation}
 \mathcal F(\zeta)=
 \begin{cases}
 \mathcal F_< (\zeta),&0\leq\zeta\leq\zeta_c,\\
 \mathcal F_< (\zeta_c)+\mathcal G(\zeta)-\mathcal G(\zeta_c),
 &\zeta_c\leq\zeta\leq\zeta_H.
 \end{cases}
 \label{eq:accepted-primitive-joined}
\end{equation}
Direct differentiation of these expressions gives
$\mathcal F'(\zeta)=Q_0(d\zeta/a)/[\Aref(1+\zeta)]$.
Therefore the closed profile is
\begin{equation}
 \boxed{k_{w,0}(z)=K+\frac{\Aref}{a}
 \left[\mathcal F\!\left(\frac{az}{d}\right)
       -\mathcal F(\zeta_H)\right].}
 \label{eq:accepted-k-closed}
\end{equation}
The integral form \cref{eq:accepted-k-profile} remains applicable
to other viscosity segments and positive diffusivities.

\paragraph{The sign of the local imbalance.}
For every $z\geq z_c>0$ on this rough branch,
\begin{equation}
 \boxed{\ECw(z)-p_{A,0}(z)
 =\frac{\Aref[(d+\nu)az+d^2]}{z(d+az)^2}>0.}
 \label{eq:accepted-tail-defect}
\end{equation}
This strict inequality concerns the $1/z$ tail. It does not imply
positive imbalance throughout the plateau. In fact, the budget
requires a compensating region of negative imbalance.

\begin{proposition}[A positive rough-wall reference]
Assume $z_v\leq0$, $d\geq2\nu$, and $\Lambda\geq1$.
Then $\ECw-p_{A,0}$ has exactly one zero in $(0,z_c)$,
changing from negative to positive. For every positive $D_{k,w}$,
\begin{equation}
 Q_0(z)<0\quad(0<z<H),\qquad
 k_{w,0}'(z)<0\quad(0<z<H),\qquad
 k_{w,0}(z)>K>0\quad(0\leq z<H).
 \label{eq:accepted-positivity}
\end{equation}
\end{proposition}
\begin{proof}
The reference production satisfies
\[
 p_{A,0}'(z)=
 \frac{a t_0^2[2\nu-(d+az)]}{(d+az)^3}\leq0,
\]
and is strictly decreasing for $z>0$. On the dissipation plateau,
$\Aref/z_c-p_{A,0}$ is therefore strictly increasing and is
positive at $z_c$ by \cref{eq:accepted-tail-defect}.
Its initial value must be negative, since otherwise the total
integral could not vanish. It crosses zero exactly once, after
which the defect remains positive up to $H$. Its cumulative
integral first decreases and then increases to $Q_0(H)=0$, so
$Q_0<0$ in the interior. The profile and monotonicity statements
follow from \cref{eq:accepted-k-profile}.
\end{proof}
The condition $d\geq2\nu$ is sufficient, not necessary.
For other intercepts, including a smooth-wall viscosity segment,
the budget construction remains defined when its gate holds, but
positivity must be assessed from the resulting profile.

The rough-branch velocity anchor also gives the useful identities
\begin{equation}
 \boxed{z_{0,c}=\frac{dH}{aH+d},\qquad
 z_c=\ee z_{0,c}
 \exp\!\left[\nu\left(\frac1d-\frac1{aH+d}\right)\right].}
 \label{eq:accepted-roughness-cutoff}
\end{equation}
Even with the calibrated slope, $z_{0,c}$ retains a finite-height
dependence. For $aH\gg d$, it tends to $d/a$.

\subsection{Conservative matching and normal pressure balance}
\label{app:accepted-awf-matching}

The inner momentum reference and the outer profile in
\cref{eq:accepted-outer-reference} satisfy
\begin{equation}
 \begin{gathered}
 U_{w,0}(0)=0,\qquad U_{w,0}(H)=U_{o,0}(H),\\
 D_A(H)U_{w,0}'(H)
 =\nu_{t,o,0}(H)U_{o,0}'(H)=t_0.
 \end{gathered}
 \label{eq:accepted-momentum-match}
\end{equation}
The calibrated coefficients match the viscosity slope, not the
finite-height viscosities: on the rough branch,
$D_A(H)=aH+d$ while $\nu_{t,o,0}(H)=aH$.
The velocity derivatives therefore need not coincide.

The turbulence matching conditions are
\begin{equation}
 \boxed{\begin{aligned}
 k_{w,0}(H)&=k_{o,0}(H)=K,\\
 D_{k,w}(H)k_{w,0}'(H)
 &=D_{k,o}(H)k_{o,0}'(H)=0,\\
 \ECw(H)&=\eps_{o,0}(H)=\frac{\Aref}{H}.
 \end{aligned}}
 \label{eq:accepted-turbulence-match}
\end{equation}
Here $D_{k,o}=\nu_t^{CLS}/\sigma_k$ on the declared outer branch.
Inner dissipation is dependent, so the outer dissipation equation
receives its single trace. No additional dissipation-flux condition
is imposed. In the limiting case $z_c=H$, even the one-sided
dissipation derivatives differ, while the required trace still
matches.

With the full-strain inner stress completion of
\cref{eq:curve-inner-stress}, the parallel inner normal stress is
$R_{zz,w}=2k_w/3$. Its normal momentum balance therefore requires
\begin{equation}
 \frac{\dd}{\dd z}\left(
 \frac{p_w}{\rho}+gz+\frac23 k_{w,0}\right)=0.
 \label{eq:accepted-inner-pressure}
\end{equation}
Mechanical pressure inside this layer is generally not hydrostatic.
Since $R_{zz,o}=0.52K$ at the outer reference, continuity of total
normal traction gives
\begin{equation}
 \boxed{\frac{p_w(H)-p_o(H)}{\rho}
 =\left(0.52-\frac23\right)K=-\frac{11}{75}K.}
 \label{eq:accepted-pressure-jump}
\end{equation}
The interface offset is unchanged from the full-strain EP-AWF
completion because the matching value is still $K$; the interior
pressure distribution changes with $k_{w,0}(z)$.
When a reduced pressure is stored, recover the mechanical pressure
consistently before evaluating the AWF pressure forcing.

The flat reference also retains the integrated mechanical-energy
identity
\begin{equation}
 t_0 U_{H,0}
 =\nu\int_0^H[U_{w,0}'(z)]^2\,\dd z
  +\int_0^H p_{A,0}(z)\,\dd z
 =\nu t_0^2J_2(H)+\int_0^H\ECw(z)\,\dd z.
 \label{eq:accepted-energy}
\end{equation}
The nonzero internal $k$ flux redistributes energy within the layer;
it contributes no net flux at the two reference endpoints.

\subsection{Distinction from EP-AWF with unit equilibrium weight}
\label{app:accepted-awf-versus-ep}

To isolate the transport difference, compare the two constructions
at identical $K,h,H,\nu$, viscosity coefficients, and budget cutoff.
At $\chi=1$, the EP-AWF has
\begin{equation}
 E_w^{EP}=\ECw,\qquad
 C_{k,EP}^w=\ECw-p_{A,0},\qquad
 S_{k,w}^{EP}=P_A-p_{A,0}.
 \label{eq:accepted-ep-comparison}
\end{equation}
At the flat reference, $P_A=p_{A,0}$ and this source vanishes
pointwise, permitting constant inner $k=K$. The accepted alternative
instead uses
\begin{equation}
 E_w=\ECw,\qquad C_k^w=0,\qquad
 S_{k,w}=P_A-\ECw.
 \label{eq:accepted-source-comparison}
\end{equation}
Its reference source is generally nonzero and is balanced by
$\partial_z(D_{k,w}k_{w,0}')$.

Both sources have the same layer integral,
\begin{equation}
 \int_0^H S_{k,w}\,\dd z
 =\int_0^H S_{k,w}^{EP}\,\dd z
 =\int_0^H(P_A-p_{A,0})\,\dd z,
 \label{eq:accepted-equal-integrated-sources}
\end{equation}
but their local balances and inner kinetic-energy fields differ.
On the reference, EP-AWF carries the physical redistribution flux
$-Q_0$ through its prescribed correction; the accepted model carries
the same flux through ordinary diffusion of its nonconstant $k$.
This reference identity does not make their evolving dynamics
equivalent. In addition, comparison with an EP-AWF that retains the
original coefficients must account separately for the changed
momentum reconstruction.

Setting the old switch to $\chi=0$ does not select the present
variant: it restores the original-cutoff dissipation $\EAw$ through
\cref{eq:blend}. The accepted model retains $\ECw$ while removing
$C_k^w$. Consequently the spatial fade to the original AWF in
\cref{sec:transition} is not part of this variant.

\subsection{A specified extension to curved terrain}
\label{app:accepted-awf-curved}

The accepted scalar policy can be used with the conservative curved
patch of \cref{sec:curved}. To make that extension definite, adopt
the note's metric-weighted reference budget, together with
$r_K=r_m=H$, $K=k|_{\Gamma_H}$, the calibrated viscosity, and
$C_k^w=0$. This is an explicit extension of the flat model; the
flat matching theorem does not by itself select a unique curved
calibration or establish an equilibrium over a hill.

Retain the plus curvature convention
\begin{equation}
 \bm x(\bm\xi,r)=\bm X(\bm\xi)+r\bm n(\bm\xi),\qquad
 \mathsf B=\nabla_\Gamma\bm n,\qquad
 \mathsf A=I_T+r\mathsf B,\qquad J=\det_T\mathsf A.
 \label{eq:accepted-curved-map}
\end{equation}
Here $r$ is physical normal distance and $\dd V=J\,\dd S_w\,\dd r$.
All principal stretches $1+r\kappa_i$ must stay positive, and
the normal map must be globally one-to-one. Positive local $J$
alone does not guarantee unique normal projection.

\paragraph{Reference cutoff in the physical volume measure.}
Define
\begin{equation}
 \boxed{\begin{aligned}
 E_w(r)&=\ECw(r)=\frac{\Aref(K)}{\max(r,r_c)},\\
 \int_0^H J(r)\ECw(r)\,\dd r
 &=\mathcal P_J:=\int_0^H J(r)p_{A,0}(r;K)\,\dd r,\\
 p_{A,0}(r;K)&=t_0(K)^2
              \frac{\nu_{t,A}(r;K)}{D_A(r;K)^2}.
 \end{aligned}}
 \label{eq:accepted-curved-budget}
\end{equation}
The reference kernel remains counterfactual flat, parallel,
zero-forcing production at the current $K$. It is not replaced
by actual terrain production.
For $\mathcal B(c)=\int_0^H J(r)\Aref/\max(r,c)\,\dd r$,
\begin{equation}
 \mathcal B'(c)=-\frac{\Aref}{c^2}\int_0^cJ(r)\,\dd r<0,
 \qquad \lim_{c\downarrow0}\mathcal B(c)=+\infty.
 \label{eq:accepted-curved-monotonicity}
\end{equation}
Therefore a unique $0<r_c\leq H$ exists exactly when
\begin{equation}
 \boxed{\mathcal P_J\geq
 \frac{\Aref}{H}\int_0^HJ(r)\,\dd r.}
 \label{eq:accepted-curved-gate}
\end{equation}
At a fixed footpoint on a three-dimensional surface,
$J=(1+\kappa_1r)(1+\kappa_2r)$. Direct integration gives
\begin{equation}
 \begin{aligned}
 \frac{\mathcal B(c)}{\Aref}
 &=1+\ln\frac Hc
 +(\kappa_1+\kappa_2)\left(H-\frac c2\right)\\
 &\quad+\kappa_1\kappa_2
       \left(\frac{H^2}{2}-\frac{c^2}{6}\right).
 \end{aligned}
 \label{eq:accepted-curved-budget-explicit}
\end{equation}
The curved cutoff thus requires one monotone scalar root solve.
The exponential formula is recovered when $J=1$.

\paragraph{Actual inner flow and transport.}
Write $\bm U=\bm u_t+w\bm n$ and
$\bm v=\mathsf A^{-1}\bm u_t$. The inner equation is
\begin{equation}
 \boxed{\begin{aligned}
 J\partial_t k
 &+\partial_r\{J(wk-D_{k,w}\partial_rk)\}\\
 &+\operatorname{div}_\Gamma\left\{
 J[\bm v k-D_{k,w}\mathsf A^{-1}\mathsf A^{-T}
                         \nabla_\Gamma k]\right\}
 =J(P_A-\ECw).
 \end{aligned}}
 \label{eq:accepted-curved-k}
\end{equation}
Use the full continuity and momentum equations
\cref{eq:curve-continuity,eq:curve-momentum}, the stated roughness
force in \cref{eq:curve-rough-force}, and the inner completion
\begin{equation}
 \begin{gathered}
 \R_A=\tfrac23k\I-\nu_{t,A}\Sm_w,\qquad
 \mathsf T_w=-\left(\frac{p_w}{\rho}+\tfrac23k\right)\I
              +D_A\Sm_w,\\
 P_A=-\R_A:\nabla\bm U
     =\tfrac12\nu_{t,A}\Sm_w:\Sm_w
 \quad\text{for }\nabla\cdot\bm U=0.
 \end{gathered}
 \label{eq:accepted-curved-stress}
\end{equation}
The isotropic stress uses local transported $k$, not sampled $K$.
For a weakly incompressible discrete field, the production
convention discussed after \cref{eq:curve-inner-stress} retains
the corresponding $-(2k/3)\nabla\cdot\bm U$ contribution when
required by the discrete stress-work balance.

The base-normal tangential traction contains curvature and normal
motion,
\begin{equation}
 \bm t_A=D_A\left[
 \partial_r\bm u_t+\mathsf A^{-T}\nabla_\Gamma w
 -\mathsf B\mathsf A^{-1}\bm u_t\right].
 \label{eq:accepted-curved-traction}
\end{equation}
Normal momentum must likewise retain curvature acceleration and
pressure balance. Merely weighting independent flat momentum
columns by $J$ does not supply these terms. General curved and
developing flow uses the coupled patch solve of \cref{sec:curved}.

The reference budget preserves the actual non-equilibrium source:
\begin{equation}
 \int_0^H J(P_A-\ECw)\,\dd r
 =\int_0^H J(P_A-p_{A,0})\,\dd r.
 \label{eq:accepted-curved-actual-budget}
\end{equation}
This quantity need not vanish; storage, tangential transport, and
interface flux balance it. Determining the cutoff from actual
$P_A$ instead of $p_{A,0}$ would impose a different energy budget.

An auxiliary frozen normal-column construction can still be written,
\begin{equation}
 \begin{aligned}
 Q_J(r)&=\int_0^rJ(s)[\ECw(s)-p_{A,0}(s)]\,\dd s,\\
 k_{\rm col}(r)&=K-\int_r^H
            \frac{Q_J(s)}{J(s)D_{k,w}(s)}\,\dd s.
 \end{aligned}
 \label{eq:accepted-curved-column}
\end{equation}
It balances only the reference source and normal diffusion.
It is not an exact terrain-flow equilibrium, since actual
production, advection, and tangential diffusion generally differ.
Its zero base-normal endpoint flux also does not establish zero
flux through an arbitrary sloping matching surface.

\paragraph{Interface conditions.}
Retain the common velocity, total traction, and mass flux in
\cref{eq:curve-momentum-interface,eq:curve-massflux}, together with
the common $k$ trace and advective--diffusive flux in
\cref{eq:curve-k-interface}. For varying $H(\bm\xi)$, these use
\begin{equation}
 \bm N_H=\frac{\bm n-\mathsf A_H^{-T}\nabla_\Gamma H}
 {\sqrt{1+|\mathsf A_H^{-T}\nabla_\Gamma H|^2}}.
 \label{eq:accepted-curved-interface-normal}
\end{equation}
Surface derivatives in the mapped equations are taken at fixed
$r$ before restriction to $H$; in particular,
$\nabla_\Gamma K=(\nabla_\Gamma k)|_H
+(\partial_rk)|_H\nabla_\Gamma H$.
The outer dissipation equation receives only
\begin{equation}
 \eps_{o,H}=\ECw(H)=\frac{\Aref(K)}{H}.
 \label{eq:accepted-curved-epsilon-trace}
\end{equation}
The full normal-traction condition replaces the flat pressure
offset over terrain. There is no independently solved inner
dissipation equation or additional prescribed dissipation flux.

\subsection{Illustrative reference and implementation implications}
\label{app:accepted-awf-example}

For a flat reference with
\begin{equation}
 \ust=0.4\units{m\,s^{-1}},\quad
 \nu=1.5\times10^{-5}\units{m^2\,s^{-1}},\quad
 h=0.1\units{m},\quad H=2\units{m},\quad z_K=z_m=H,
 \label{eq:accepted-example-inputs}
\end{equation}
the unchanged roughness correlation and calibrated coefficients give
the values in \cref{tab:accepted-reference}. They evaluate the
analytical construction; they are not a solver-validation result.

\begin{table}[htbp]
\centering\small
\begin{tabular}{@{}lll@{}}
\toprule
Quantity & Value & Units\\
\midrule
$K$ & $0.5105080131$ & $\mathrm{m^2\,s^{-2}}$\\
$z_v$ & $-0.0052721312$ & $\mathrm m$\\
$a$ & $0.1762215621$ & $\mathrm{m\,s^{-1}}$\\
$d$ & $0.0009440632$ & $\mathrm{m^2\,s^{-1}}$\\
$\mathcal P_0$ & $0.8584509755$ & $\mathrm{m^3\,s^{-3}}$\\
$z_c$ & $0.0147555946$ & $\mathrm m$\\
$U_{H,0}$ & $5.3797062304$ & $\mathrm{m\,s^{-1}}$\\
$z_{0,c}$ & $0.0053429396$ & $\mathrm m$\\
$k_{w,0}(0)$ & $1.0229104395$ & $\mathrm{m^2\,s^{-2}}$\\
$k_{w,0}(0.1\units m)$ & $0.5447982810$ & $\mathrm{m^2\,s^{-2}}$\\
$k_{w,0}(1\units m)$ & $0.5111909311$ & $\mathrm{m^2\,s^{-2}}$\\
\bottomrule
\end{tabular}
\caption{Accepted flat-wall reference with endpoint sampling.
The outer kinetic energy remains $K$, while ordinary diffusion
supports the larger inner values.}
\label{tab:accepted-reference}
\end{table}

Here $d>2\nu$, so the positivity proposition applies. In particular,
$\ECw(0)\simeq9.84519\units{m^2\,s^{-3}}$ is smaller than
$p_{A,0}(0)\simeq26.68598\units{m^2\,s^{-3}}$.
The defect changes sign near $z=0.00931273\units m$ and remains
positive thereafter, while its full-layer integral is zero.
These signs refer to the accepted budget cutoff; retaining the
original molecular cutoff would give a different imbalance.

Implementation should evaluate $K$ from the common matching trace,
construct $D_A$ with the calibrated coefficients, and determine
the cutoff from reference production at that same $K$.
Actual production is used in \cref{eq:accepted-k-transport} or
\cref{eq:accepted-curved-k}. Initialize a flat reference with
\cref{eq:accepted-k-profile}, retain zero modeled ground $k$ flux,
and couple the matching value and conservative flux to outer CLS.
Do not reintroduce the EP redistribution while using this profile.
Quadrature should resolve each active viscosity, forcing, and
dissipation breakpoint, and the full assembled reference residual
must be checked with the actual interface and pressure operators.

The derivation establishes continuum compatibility with the
standard outer CLS equilibrium on flat, undisturbed terrain and
specifies a conservative curved extension. It does not require
the outer flow over a hill to remain logarithmic. The positive
flat reference is also not a positivity or stability theorem for
the evolving equations: a sink controlled by sampled $K$ can
remain positive as local $k_w$ decreases, and the constitutive
limitations of the unchanged outer CLS model identified in
\cref{sec:stability} remain. Any active robustness modification
must be assessed against this reference and declared separately.

% The main note already calls \appendix. Do not call it again here.
% Uses the main note's macros, theorem environments, labels and bibliography.

\section{AWF with CLS-compatible coefficients, a budget-determined cutoff,
and a nonconstant inner kinetic-energy profile}
\label{app:accepted-awf}

This appendix defines the accepted alternative to the constant-inner-$k$
EP-AWF construction. The objective is to preserve the standard neutral
CLS equilibrium in the outer region while allowing the kinetic energy
to vary inside the modeled wall layer. The AWF viscosity coefficients
are calibrated to that outer equilibrium, dissipation retains the
budget-determined cutoff, and the prescribed redistribution $C_k^w$
is removed. Ordinary diffusion of the transported inner $k$ then
balances the local difference between production and dissipation.
The outer CLS model, including its Yap correction, is unchanged.

The original CLS and rough-wall AWF supply the constitutive and
wall-model starting points \cite{CLS1996,SCI2006}. The coefficient
calibration, budget cutoff, retained inner transport, and curved
completion specified here are modifications of that combination.
When this variant is selected, the definitions below replace the
original wall coefficients, interior sampling choice, and
$\chi$-dependent source policy used elsewhere in the note.

\subsection{Scope, endpoint sampling, and coefficient calibration}
\label{app:accepted-awf-coefficients}

Use the high-Re outer equation convention of
\cref{sec:derivation}: molecular momentum and scalar diffusion are
omitted in the outer analytical reference, while $\nu>0$ is retained
in the AWF. This is a statement about the selected equations, not an
assertion that the same profiles solve a finite-Re outer system.
The transported inner turbulence variable is $k$; inner
dissipation is a dependent total $k$ sink. No inner CLS dissipation
transport equation is introduced.

The accepted sampling and tail-matching heights are
\begin{equation}
 \boxed{z_K=z_m=H,\qquad
 K=k_w(H)=k_o(H)>0.}
 \label{eq:accepted-sampling}
\end{equation}
The subscripts $w,o$ distinguish the modeled wall layer and outer
region. At each positive current $K$, define the counterfactual
reference scales
\begin{equation}
 \ust(K)=\sqrt{\frac{K}{\qs}},\qquad
 t_0(K)=\frac{K}{\qs},\qquad
 \Aref(K)=\frac{\ust(K)^3}{\kc}.
 \label{eq:accepted-reference-scales}
\end{equation}
These are reference scales. In a developing, forced, or separated
flow, $t_0(K)$ is not imposed as the actual wall traction.
Arguments $K$ are suppressed below when all quantities use the same
sample. A fixed inlet reference is obtained by setting $K=K_0$.

The outer constants are the roots derived in
\cref{eq:qroot,eq:kappaeq},
\begin{equation}
 \begin{gathered}
 C(\qs)\qs^2=1,\qquad
 \frac{\qs\kc^2}{1.3}
 +H_Y\!\left(\frac{\kc\qs^{3/2}}{2.5}\right)=0.48,\\
 \qs=3.190675081624288,\qquad
 \kc=0.440553905186659,\qquad
 \Cs=\qs^{-2}=0.0982278953282113.
 \end{gathered}
 \label{eq:accepted-cls-constants}
\end{equation}
Here $C$ and $H_Y$ retain the definitions of the original high-Re CLS
closure and its Yap correction. Matching the dissipation coefficient
and the asymptotic auxiliary-viscosity slope gives
\begin{equation}
 \boxed{\begin{aligned}
 c_\ell^w&=\kc\qs^{3/2}=2.510862475219797,\\
 c_\mu^w&=\Cs=\qs^{-2}=0.0982278953282113,\\
 \alpha&=c_\mu^wc_\ell^w
       =\frac{\kc}{\sqrt{\qs}}=0.246636736399424.
 \end{aligned}}
 \label{eq:accepted-coefficients}
\end{equation}
Indeed, with $a:=\alpha\sqrt K$,
\begin{equation}
 \boxed{a=\kc\ust,\qquad
 \frac{K^{3/2}}{c_\ell^w}=\Aref
 =\frac{t_0^2}{a}.}
 \label{eq:accepted-calibration-identities}
\end{equation}
The choice of $c_\ell^w$ matches the outer dissipation tail exactly.
The additional choice of $c_\mu^w$ matches the viscosity slope;
conservative interface matching by itself does not require equality
of these slopes.

Retain the original intercept correlation $b(K,h,\nu)$ in
\cref{eq:brough}, including $b=10.7$ for $h=0$. Define
\begin{equation}
 z_v=\frac{\nu b}{\sqrt K},\qquad
 \nu_{t,A}(z)=a\max(z-z_v,0),\qquad
 D_A=\nu+\nu_{t,A},\qquad
 D_{k,w}=\nu+\frac{\nu_{t,A}}{\sigma_k}.
 \label{eq:accepted-viscosities}
\end{equation}
The CLS choice $\sigma_k=1$ gives $D_{k,w}=D_A$.
These auxiliary coefficients depend on sampled $K$, while the
transported field $k_w(z)$ can vary. Substituting $k_w(z)$ for $K$
in the viscosity or dissipation law would define a different model.

\subsection{Momentum, production, and the outer velocity anchor}
\label{app:accepted-awf-momentum}

The analytical form of the AWF momentum reconstruction is retained,
with all integrals evaluated using the calibrated $D_A$. Put
$L_h(z)=\max(z-h,0)$ and
\begin{equation}
 J_0(z)=\int_0^z\frac{\dd s}{D_A(s)},\qquad
 J_1(z)=\int_0^z\frac{L_h(s)}{D_A(s)}\,\dd s,\qquad
 J_2(z)=\int_0^z\frac{\dd s}{D_A(s)^2}.
 \label{eq:accepted-J}
\end{equation}
For scalar tangential forcing $F$, the flat quasi-steady reduction is
\begin{equation}
 \begin{aligned}
 t_A&=\frac{U_H-FJ_1(H)}{J_0(H)},\qquad
 U_A(z)=t_AJ_0(z)+FJ_1(z),\\
 P_A(z)&=\nu_{t,A}(z)
       \left[\frac{t_A+FL_h(z)}{D_A(z)}\right]^2.
 \end{aligned}
 \label{eq:accepted-forced-awf}
\end{equation}
Thus no slip, the original forcing cutoff at $h$, and the explicit
pressure-gradient response retain their analytical form. Their
numerical values change with the new viscosity coefficients.
The vector form follows \cref{eq:3d}.

At the counterfactual zero-forcing reference, choose
\begin{equation}
 U_{H,0}=t_0J_0(H),\qquad
 U_{w,0}(z)=t_0J_0(z),\qquad
 p_{A,0}(z)=t_0^2\frac{\nu_{t,A}(z)}{D_A(z)^2}.
 \label{eq:accepted-reference-production}
\end{equation}
The resulting integrated reference production is
\begin{equation}
 \boxed{\mathcal P_0:=\int_0^H p_{A,0}(z)\,\dd z
 =t_0^2[J_0(H)-\nu J_2(H)]=HP_0.}
 \label{eq:accepted-production-budget}
\end{equation}
$P_0$ remains the height mean used in the main note; $p_{A,0}$ is
pointwise reference production. Neither is the production of a
general current terrain flow.

For $z\geq H$, preserve the standard outer CLS branch
\begin{equation}
 \boxed{\begin{aligned}
 U_{o,0}(z)&=t_0J_0(H)+\frac{\ust}{\kc}\ln\frac zH,\\
 k_{o,0}(z)&=K,\qquad
 \eps_{o,0}(z)=\frac{\Aref}{z},\qquad
 \nu_{t,o,0}(z)=\kc\ust z.
 \end{aligned}}
 \label{eq:accepted-outer-reference}
\end{equation}
The outer velocity offset has been used to satisfy the AWF anchor.
Its aerodynamic intercept is therefore
\begin{equation}
 z_{0,c}=H\exp[-\kc\ust J_0(H)].
 \label{eq:accepted-aerodynamic-intercept}
\end{equation}
Retaining the same sand-grain $h$ while changing the AWF coefficients
generally changes this intercept. A separately prescribed $z_{0,c}$
requires an additional momentum/intercept calibration.

\subsection{Budget-determined dissipation and admissibility}
\label{app:accepted-awf-cutoff}

Use the continuous positive dissipation distribution
\begin{equation}
 \boxed{E_w(z)=\ECw(z)=\frac{\Aref}{\max(z,z_c)},
 \qquad C_k^w(z)=0.}
 \label{eq:accepted-dissipation}
\end{equation}
The required tail begins no later than $H$, so $0<z_c\leq H$.
Its integral is
\begin{equation}
 \int_0^H\ECw(z)\,\dd z
 =\Aref\left(1+\ln\frac H{z_c}\right).
 \label{eq:accepted-dissipation-integral}
\end{equation}
The standard outer reference has zero $k$ gradient. Together with
zero modeled ground $k$ flux, conservative matching therefore
requires
\begin{equation}
 \boxed{\int_0^H\ECw\,\dd z=\mathcal P_0.}
 \label{eq:accepted-budget-condition}
\end{equation}
Combining the preceding identities gives
\begin{equation}
 \boxed{z_c=H\exp(1-\Lambda),\qquad
 \Lambda=\frac{\mathcal P_0}{\Aref}
 =\frac{t_0^2[J_0(H)-\nu J_2(H)]}{\Aref}.}
 \label{eq:accepted-cutoff}
\end{equation}
The left side of \cref{eq:accepted-dissipation-integral} decreases
strictly with $z_c$, tends to infinity as $z_c\downarrow0$, and equals
$\Aref$ at $z_c=H$. Consequently,
\begin{equation}
 \boxed{0<z_c\leq H\quad\Longleftrightarrow\quad\Lambda\geq1,}
 \label{eq:accepted-cutoff-gate}
\end{equation}
and the admissible cutoff is unique. This is the $z_m=H$
specialization of \cref{eq:hp-gate}. If the gate fails, this
monotone cutoff family cannot meet both the budget and outer trace.
Clipping $z_c$ to $H$ would not preserve the budget.

The original molecular cutoff
$z_{\eps,A}=2c_\ell^w\nu/\sqrt K$ is replaced by
\cref{eq:accepted-cutoff}. The plateau is now $\Aref/z_c$.
It cannot simultaneously remain $2\nu K/z_c^2$: imposing both
plateau formulas would force the original molecular cutoff again.
Thus the cutoff change is an explicit change of the inner
dissipation model.

\subsection{Derivation of the nonconstant inner reference}
\label{app:accepted-awf-inner-k}

Retain the inner transport equation with its physical production
and prescribed dissipation,
\begin{equation}
 \partial_t k_w+\bm U_w\cdot\nabla k_w
 =\nabla\cdot(D_{k,w}\nabla k_w)+P_A-\ECw.
 \label{eq:accepted-k-transport}
\end{equation}
For the steady, parallel flat reference this reduces to
\begin{equation}
 (D_{k,w}k_{w,0}')'=\ECw-p_{A,0},\qquad
 D_{k,w}(0)k_{w,0}'(0)=0,\qquad k_{w,0}(H)=K.
 \label{eq:accepted-k-bvp}
\end{equation}
Define the cumulative source defect
\begin{equation}
 Q_0(z)=\int_0^z[\ECw(s)-p_{A,0}(s)]\,\dd s.
 \label{eq:accepted-Q}
\end{equation}
Two integrations give the unique solution for the prescribed
positive diffusivity,
\begin{equation}
 \boxed{D_{k,w}(z)k_{w,0}'(z)=Q_0(z),\qquad
 k_{w,0}(z)=K-\int_z^H\frac{Q_0(s)}{D_{k,w}(s)}\,\dd s.}
 \label{eq:accepted-k-profile}
\end{equation}
The quantity $Q_0$ is the gradient flux inside the diffusion
operator; the physical upward diffusive flux is $-Q_0$.
The cutoff budget gives
\begin{equation}
 Q_0(0)=Q_0(H)=0,\qquad k_{w,0}'(H)=0.
 \label{eq:accepted-endpoint-fluxes}
\end{equation}
Conversely, these two zero endpoint fluxes imply
\cref{eq:accepted-budget-condition} by integration of the $k$
equation. Allowing a variable $k_w$ supplies local diffusion; it
does not remove the integrated compatibility requirement.

No additional value of $k_w(0)$ is prescribed. It follows from
\cref{eq:accepted-k-profile}. This is a modeled rough-wall layer,
not a viscous-sublayer resolution with an independently imposed
$k(0)=0$. The formula defines a stationary reference and an
initialization; during a general flow, $k_w$ is evolved by
\cref{eq:accepted-k-transport}, rather than reset to this profile.

For every sufficiently regular test function $\phi$, integration
by parts yields
\begin{equation}
 \boxed{-\int_0^H D_{k,w}k_{w,0}'\phi'\,\dd z
 +\int_0^H(p_{A,0}-\ECw)\phi\,\dd z=0.}
 \label{eq:accepted-weak-equilibrium}
\end{equation}
Thus the nonconstant reference balances every continuum test
moment. The higher-mode obstruction in \cref{sec:hp-source}
assumed constant inner $k$, for which the diffusion term vanished.
Finite-dimensional projection and numerical differentiation can
still leave a discrete residual; continuum balance alone does not
establish exact spectral/$hp$ preservation.

\subsection{Closed formulas and positivity on the rough branch}
\label{app:accepted-awf-rough}

Suppose $z_v\leq0$ and define
\begin{equation}
 d=\nu-az_v\geq\nu,\qquad D_A(z)=d+az.
 \label{eq:accepted-rough-D}
\end{equation}
Then
\begin{equation}
 \begin{aligned}
 J_0(z)&=\frac1a\ln\frac{d+az}{d},&
 J_2(z)&=\frac1a\left(\frac1d-\frac1{d+az}\right),\\
 p_{A,0}(z)&=t_0^2\frac{d+az-\nu}{(d+az)^2}.&&
 \end{aligned}
 \label{eq:accepted-rough-integrals}
\end{equation}
Using $\Aref=t_0^2/a$, the cumulative production is
\begin{equation}
 \int_0^z p_{A,0}(s)\,\dd s
 =\Aref\left[
 \ln\frac{d+az}{d}
 -\nu\left(\frac1d-\frac1{d+az}\right)\right].
 \label{eq:accepted-rough-production-integral}
\end{equation}
In particular, with $\zeta_H=aH/d$ and $\beta=\nu/d$,
\begin{equation}
 \Lambda=\ln(1+\zeta_H)
          -\frac{\beta\zeta_H}{1+\zeta_H}.
 \label{eq:accepted-rough-Lambda}
\end{equation}
The cumulative dissipation is $\Aref R_c(z)$, where
\begin{equation}
 R_c(z)=\begin{cases}
 z/z_c,&0\leq z\leq z_c,\\
 1+\ln(z/z_c),&z_c\leq z\leq H.
 \end{cases}
 \label{eq:accepted-Rc}
\end{equation}
Hence an elementary expression for the flux is
\begin{equation}
 Q_0(z)=\Aref\left[
 R_c(z)-\ln\frac{d+az}{d}
 +\nu\left(\frac1d-\frac1{d+az}\right)\right].
 \label{eq:accepted-rough-Q}
\end{equation}
Together with \cref{eq:accepted-k-profile}, this provides an exact
analytical representation of the inner reference.

\paragraph{A closed primitive for the kinetic energy.}
For the selected $\sigma_k=1$, set
$\zeta=az/d$, $\zeta_c=az_c/d$, and define
\begin{align}
 \mathcal F_< (\zeta)
 &=\frac{\zeta-\ln(1+\zeta)}{\zeta_c}
   -\frac12\ln^2(1+\zeta)
   +\beta\left[\ln(1+\zeta)+\frac1{1+\zeta}-1\right],
 \label{eq:accepted-primitive-lower}\\
 \mathcal G(\zeta)
 &=\left[1+\ln\frac{\zeta}{\zeta_c}+\beta\right]\ln(1+\zeta)
   +\operatorname{Li}_2(-\zeta)
   -\frac12\ln^2(1+\zeta)+\frac{\beta}{1+\zeta},
 \quad \zeta>0.
 \label{eq:accepted-primitive-upper}
\end{align}
The real dilogarithm is defined here by
$\operatorname{Li}_2(-\zeta)=-\int_0^\zeta\ln(1+s)\,\dd s/s$.
Join the primitives continuously:
\begin{equation}
 \mathcal F(\zeta)=
 \begin{cases}
 \mathcal F_< (\zeta),&0\leq\zeta\leq\zeta_c,\\
 \mathcal F_< (\zeta_c)+\mathcal G(\zeta)-\mathcal G(\zeta_c),
 &\zeta_c\leq\zeta\leq\zeta_H.
 \end{cases}
 \label{eq:accepted-primitive-joined}
\end{equation}
Direct differentiation of these expressions gives
$\mathcal F'(\zeta)=Q_0(d\zeta/a)/[\Aref(1+\zeta)]$.
Therefore the closed profile is
\begin{equation}
 \boxed{k_{w,0}(z)=K+\frac{\Aref}{a}
 \left[\mathcal F\!\left(\frac{az}{d}\right)
       -\mathcal F(\zeta_H)\right].}
 \label{eq:accepted-k-closed}
\end{equation}
The integral form \cref{eq:accepted-k-profile} remains applicable
to other viscosity segments and positive diffusivities.

\paragraph{The sign of the local imbalance.}
For every $z\geq z_c>0$ on this rough branch,
\begin{equation}
 \boxed{\ECw(z)-p_{A,0}(z)
 =\frac{\Aref[(d+\nu)az+d^2]}{z(d+az)^2}>0.}
 \label{eq:accepted-tail-defect}
\end{equation}
This strict inequality concerns the $1/z$ tail. It does not imply
positive imbalance throughout the plateau. In fact, the budget
requires a compensating region of negative imbalance.

\begin{proposition}[A positive rough-wall reference]
Assume $z_v\leq0$, $d\geq2\nu$, and $\Lambda\geq1$.
Then $\ECw-p_{A,0}$ has exactly one zero in $(0,z_c)$,
changing from negative to positive. For every positive $D_{k,w}$,
\begin{equation}
 Q_0(z)<0\quad(0<z<H),\qquad
 k_{w,0}'(z)<0\quad(0<z<H),\qquad
 k_{w,0}(z)>K>0\quad(0\leq z<H).
 \label{eq:accepted-positivity}
\end{equation}
\end{proposition}
\begin{proof}
The reference production satisfies
\[
 p_{A,0}'(z)=
 \frac{a t_0^2[2\nu-(d+az)]}{(d+az)^3}\leq0,
\]
and is strictly decreasing for $z>0$. On the dissipation plateau,
$\Aref/z_c-p_{A,0}$ is therefore strictly increasing and is
positive at $z_c$ by \cref{eq:accepted-tail-defect}.
Its initial value must be negative, since otherwise the total
integral could not vanish. It crosses zero exactly once, after
which the defect remains positive up to $H$. Its cumulative
integral first decreases and then increases to $Q_0(H)=0$, so
$Q_0<0$ in the interior. The profile and monotonicity statements
follow from \cref{eq:accepted-k-profile}.
\end{proof}
The condition $d\geq2\nu$ is sufficient, not necessary.
For other intercepts, including a smooth-wall viscosity segment,
the budget construction remains defined when its gate holds, but
positivity must be assessed from the resulting profile.

The rough-branch velocity anchor also gives the useful identities
\begin{equation}
 \boxed{z_{0,c}=\frac{dH}{aH+d},\qquad
 z_c=\ee z_{0,c}
 \exp\!\left[\nu\left(\frac1d-\frac1{aH+d}\right)\right].}
 \label{eq:accepted-roughness-cutoff}
\end{equation}
Even with the calibrated slope, $z_{0,c}$ retains a finite-height
dependence. For $aH\gg d$, it tends to $d/a$.

\subsection{Conservative matching and normal pressure balance}
\label{app:accepted-awf-matching}

The inner momentum reference and the outer profile in
\cref{eq:accepted-outer-reference} satisfy
\begin{equation}
 \begin{gathered}
 U_{w,0}(0)=0,\qquad U_{w,0}(H)=U_{o,0}(H),\\
 D_A(H)U_{w,0}'(H)
 =\nu_{t,o,0}(H)U_{o,0}'(H)=t_0.
 \end{gathered}
 \label{eq:accepted-momentum-match}
\end{equation}
The calibrated coefficients match the viscosity slope, not the
finite-height viscosities: on the rough branch,
$D_A(H)=aH+d$ while $\nu_{t,o,0}(H)=aH$.
The velocity derivatives therefore need not coincide.

The turbulence matching conditions are
\begin{equation}
 \boxed{\begin{aligned}
 k_{w,0}(H)&=k_{o,0}(H)=K,\\
 D_{k,w}(H)k_{w,0}'(H)
 &=D_{k,o}(H)k_{o,0}'(H)=0,\\
 \ECw(H)&=\eps_{o,0}(H)=\frac{\Aref}{H}.
 \end{aligned}}
 \label{eq:accepted-turbulence-match}
\end{equation}
Here $D_{k,o}=\nu_t^{CLS}/\sigma_k$ on the declared outer branch.
Inner dissipation is dependent, so the outer dissipation equation
receives its single trace. No additional dissipation-flux condition
is imposed. In the limiting case $z_c=H$, even the one-sided
dissipation derivatives differ, while the required trace still
matches.

With the full-strain inner stress completion of
\cref{eq:curve-inner-stress}, the parallel inner normal stress is
$R_{zz,w}=2k_w/3$. Its normal momentum balance therefore requires
\begin{equation}
 \frac{\dd}{\dd z}\left(
 \frac{p_w}{\rho}+gz+\frac23 k_{w,0}\right)=0.
 \label{eq:accepted-inner-pressure}
\end{equation}
Mechanical pressure inside this layer is generally not hydrostatic.
Since $R_{zz,o}=0.52K$ at the outer reference, continuity of total
normal traction gives
\begin{equation}
 \boxed{\frac{p_w(H)-p_o(H)}{\rho}
 =\left(0.52-\frac23\right)K=-\frac{11}{75}K.}
 \label{eq:accepted-pressure-jump}
\end{equation}
The interface offset is unchanged from the full-strain EP-AWF
completion because the matching value is still $K$; the interior
pressure distribution changes with $k_{w,0}(z)$.
When a reduced pressure is stored, recover the mechanical pressure
consistently before evaluating the AWF pressure forcing.

The flat reference also retains the integrated mechanical-energy
identity
\begin{equation}
 t_0 U_{H,0}
 =\nu\int_0^H[U_{w,0}'(z)]^2\,\dd z
  +\int_0^H p_{A,0}(z)\,\dd z
 =\nu t_0^2J_2(H)+\int_0^H\ECw(z)\,\dd z.
 \label{eq:accepted-energy}
\end{equation}
The nonzero internal $k$ flux redistributes energy within the layer;
it contributes no net flux at the two reference endpoints.

\subsection{Distinction from EP-AWF with unit equilibrium weight}
\label{app:accepted-awf-versus-ep}

To isolate the transport difference, compare the two constructions
at identical $K,h,H,\nu$, viscosity coefficients, and budget cutoff.
At $\chi=1$, the EP-AWF has
\begin{equation}
 E_w^{EP}=\ECw,\qquad
 C_{k,EP}^w=\ECw-p_{A,0},\qquad
 S_{k,w}^{EP}=P_A-p_{A,0}.
 \label{eq:accepted-ep-comparison}
\end{equation}
At the flat reference, $P_A=p_{A,0}$ and this source vanishes
pointwise, permitting constant inner $k=K$. The accepted alternative
instead uses
\begin{equation}
 E_w=\ECw,\qquad C_k^w=0,\qquad
 S_{k,w}=P_A-\ECw.
 \label{eq:accepted-source-comparison}
\end{equation}
Its reference source is generally nonzero and is balanced by
$\partial_z(D_{k,w}k_{w,0}')$.

Both sources have the same layer integral,
\begin{equation}
 \int_0^H S_{k,w}\,\dd z
 =\int_0^H S_{k,w}^{EP}\,\dd z
 =\int_0^H(P_A-p_{A,0})\,\dd z,
 \label{eq:accepted-equal-integrated-sources}
\end{equation}
but their local balances and inner kinetic-energy fields differ.
On the reference, EP-AWF carries the physical redistribution flux
$-Q_0$ through its prescribed correction; the accepted model carries
the same flux through ordinary diffusion of its nonconstant $k$.
This reference identity does not make their evolving dynamics
equivalent. In addition, comparison with an EP-AWF that retains the
original coefficients must account separately for the changed
momentum reconstruction.

Setting the old switch to $\chi=0$ does not select the present
variant: it restores the original-cutoff dissipation $\EAw$ through
\cref{eq:blend}. The accepted model retains $\ECw$ while removing
$C_k^w$. Consequently the spatial fade to the original AWF in
\cref{sec:transition} is not part of this variant.

\subsection{A specified extension to curved terrain}
\label{app:accepted-awf-curved}

The accepted scalar policy can be used with the conservative curved
patch of \cref{sec:curved}. To make that extension definite, adopt
the note's metric-weighted reference budget, together with
$r_K=r_m=H$, $K=k|_{\Gamma_H}$, the calibrated viscosity, and
$C_k^w=0$. This is an explicit extension of the flat model; the
flat matching theorem does not by itself select a unique curved
calibration or establish an equilibrium over a hill.

Retain the plus curvature convention
\begin{equation}
 \bm x(\bm\xi,r)=\bm X(\bm\xi)+r\bm n(\bm\xi),\qquad
 \mathsf B=\nabla_\Gamma\bm n,\qquad
 \mathsf A=I_T+r\mathsf B,\qquad J=\det_T\mathsf A.
 \label{eq:accepted-curved-map}
\end{equation}
Here $r$ is physical normal distance and $\dd V=J\,\dd S_w\,\dd r$.
All principal stretches $1+r\kappa_i$ must stay positive, and
the normal map must be globally one-to-one. Positive local $J$
alone does not guarantee unique normal projection.

\paragraph{Reference cutoff in the physical volume measure.}
Define
\begin{equation}
 \boxed{\begin{aligned}
 E_w(r)&=\ECw(r)=\frac{\Aref(K)}{\max(r,r_c)},\\
 \int_0^H J(r)\ECw(r)\,\dd r
 &=\mathcal P_J:=\int_0^H J(r)p_{A,0}(r;K)\,\dd r,\\
 p_{A,0}(r;K)&=t_0(K)^2
              \frac{\nu_{t,A}(r;K)}{D_A(r;K)^2}.
 \end{aligned}}
 \label{eq:accepted-curved-budget}
\end{equation}
The reference kernel remains counterfactual flat, parallel,
zero-forcing production at the current $K$. It is not replaced
by actual terrain production.
For $\mathcal B(c)=\int_0^H J(r)\Aref/\max(r,c)\,\dd r$,
\begin{equation}
 \mathcal B'(c)=-\frac{\Aref}{c^2}\int_0^cJ(r)\,\dd r<0,
 \qquad \lim_{c\downarrow0}\mathcal B(c)=+\infty.
 \label{eq:accepted-curved-monotonicity}
\end{equation}
Therefore a unique $0<r_c\leq H$ exists exactly when
\begin{equation}
 \boxed{\mathcal P_J\geq
 \frac{\Aref}{H}\int_0^HJ(r)\,\dd r.}
 \label{eq:accepted-curved-gate}
\end{equation}
At a fixed footpoint on a three-dimensional surface,
$J=(1+\kappa_1r)(1+\kappa_2r)$. Direct integration gives
\begin{equation}
 \begin{aligned}
 \frac{\mathcal B(c)}{\Aref}
 &=1+\ln\frac Hc
 +(\kappa_1+\kappa_2)\left(H-\frac c2\right)\\
 &\quad+\kappa_1\kappa_2
       \left(\frac{H^2}{2}-\frac{c^2}{6}\right).
 \end{aligned}
 \label{eq:accepted-curved-budget-explicit}
\end{equation}
The curved cutoff thus requires one monotone scalar root solve.
The exponential formula is recovered when $J=1$.

\paragraph{Actual inner flow and transport.}
Write $\bm U=\bm u_t+w\bm n$ and
$\bm v=\mathsf A^{-1}\bm u_t$. The inner equation is
\begin{equation}
 \boxed{\begin{aligned}
 J\partial_t k
 &+\partial_r\{J(wk-D_{k,w}\partial_rk)\}\\
 &+\operatorname{div}_\Gamma\left\{
 J[\bm v k-D_{k,w}\mathsf A^{-1}\mathsf A^{-T}
                         \nabla_\Gamma k]\right\}
 =J(P_A-\ECw).
 \end{aligned}}
 \label{eq:accepted-curved-k}
\end{equation}
Use the full continuity and momentum equations
\cref{eq:curve-continuity,eq:curve-momentum}, the stated roughness
force in \cref{eq:curve-rough-force}, and the inner completion
\begin{equation}
 \begin{gathered}
 \R_A=\tfrac23k\I-\nu_{t,A}\Sm_w,\qquad
 \mathsf T_w=-\left(\frac{p_w}{\rho}+\tfrac23k\right)\I
              +D_A\Sm_w,\\
 P_A=-\R_A:\nabla\bm U
     =\tfrac12\nu_{t,A}\Sm_w:\Sm_w
 \quad\text{for }\nabla\cdot\bm U=0.
 \end{gathered}
 \label{eq:accepted-curved-stress}
\end{equation}
The isotropic stress uses local transported $k$, not sampled $K$.
For a weakly incompressible discrete field, the production
convention discussed after \cref{eq:curve-inner-stress} retains
the corresponding $-(2k/3)\nabla\cdot\bm U$ contribution when
required by the discrete stress-work balance.

The base-normal tangential traction contains curvature and normal
motion,
\begin{equation}
 \bm t_A=D_A\left[
 \partial_r\bm u_t+\mathsf A^{-T}\nabla_\Gamma w
 -\mathsf B\mathsf A^{-1}\bm u_t\right].
 \label{eq:accepted-curved-traction}
\end{equation}
Normal momentum must likewise retain curvature acceleration and
pressure balance. Merely weighting independent flat momentum
columns by $J$ does not supply these terms. General curved and
developing flow uses the coupled patch solve of \cref{sec:curved}.

The reference budget preserves the actual non-equilibrium source:
\begin{equation}
 \int_0^H J(P_A-\ECw)\,\dd r
 =\int_0^H J(P_A-p_{A,0})\,\dd r.
 \label{eq:accepted-curved-actual-budget}
\end{equation}
This quantity need not vanish; storage, tangential transport, and
interface flux balance it. Determining the cutoff from actual
$P_A$ instead of $p_{A,0}$ would impose a different energy budget.

An auxiliary frozen normal-column construction can still be written,
\begin{equation}
 \begin{aligned}
 Q_J(r)&=\int_0^rJ(s)[\ECw(s)-p_{A,0}(s)]\,\dd s,\\
 k_{\rm col}(r)&=K-\int_r^H
            \frac{Q_J(s)}{J(s)D_{k,w}(s)}\,\dd s.
 \end{aligned}
 \label{eq:accepted-curved-column}
\end{equation}
It balances only the reference source and normal diffusion.
It is not an exact terrain-flow equilibrium, since actual
production, advection, and tangential diffusion generally differ.
Its zero base-normal endpoint flux also does not establish zero
flux through an arbitrary sloping matching surface.

\paragraph{Interface conditions.}
Retain the common velocity, total traction, and mass flux in
\cref{eq:curve-momentum-interface,eq:curve-massflux}, together with
the common $k$ trace and advective--diffusive flux in
\cref{eq:curve-k-interface}. For varying $H(\bm\xi)$, these use
\begin{equation}
 \bm N_H=\frac{\bm n-\mathsf A_H^{-T}\nabla_\Gamma H}
 {\sqrt{1+|\mathsf A_H^{-T}\nabla_\Gamma H|^2}}.
 \label{eq:accepted-curved-interface-normal}
\end{equation}
Surface derivatives in the mapped equations are taken at fixed
$r$ before restriction to $H$; in particular,
$\nabla_\Gamma K=(\nabla_\Gamma k)|_H
+(\partial_rk)|_H\nabla_\Gamma H$.
The outer dissipation equation receives only
\begin{equation}
 \eps_{o,H}=\ECw(H)=\frac{\Aref(K)}{H}.
 \label{eq:accepted-curved-epsilon-trace}
\end{equation}
The full normal-traction condition replaces the flat pressure
offset over terrain. There is no independently solved inner
dissipation equation or additional prescribed dissipation flux.

\subsection{Illustrative reference and implementation implications}
\label{app:accepted-awf-example}

For a flat reference with
\begin{equation}
 \ust=0.4\units{m\,s^{-1}},\quad
 \nu=1.5\times10^{-5}\units{m^2\,s^{-1}},\quad
 h=0.1\units{m},\quad H=2\units{m},\quad z_K=z_m=H,
 \label{eq:accepted-example-inputs}
\end{equation}
the unchanged roughness correlation and calibrated coefficients give
the values in \cref{tab:accepted-reference}. They evaluate the
analytical construction; they are not a solver-validation result.

\begin{table}[htbp]
\centering\small
\begin{tabular}{@{}lll@{}}
\toprule
Quantity & Value & Units\\
\midrule
$K$ & $0.5105080131$ & $\mathrm{m^2\,s^{-2}}$\\
$z_v$ & $-0.0052721312$ & $\mathrm m$\\
$a$ & $0.1762215621$ & $\mathrm{m\,s^{-1}}$\\
$d$ & $0.0009440632$ & $\mathrm{m^2\,s^{-1}}$\\
$\mathcal P_0$ & $0.8584509755$ & $\mathrm{m^3\,s^{-3}}$\\
$z_c$ & $0.0147555946$ & $\mathrm m$\\
$U_{H,0}$ & $5.3797062304$ & $\mathrm{m\,s^{-1}}$\\
$z_{0,c}$ & $0.0053429396$ & $\mathrm m$\\
$k_{w,0}(0)$ & $1.0229104395$ & $\mathrm{m^2\,s^{-2}}$\\
$k_{w,0}(0.1\units m)$ & $0.5447982810$ & $\mathrm{m^2\,s^{-2}}$\\
$k_{w,0}(1\units m)$ & $0.5111909311$ & $\mathrm{m^2\,s^{-2}}$\\
\bottomrule
\end{tabular}
\caption{Accepted flat-wall reference with endpoint sampling.
The outer kinetic energy remains $K$, while ordinary diffusion
supports the larger inner values.}
\label{tab:accepted-reference}
\end{table}

Here $d>2\nu$, so the positivity proposition applies. In particular,
$\ECw(0)\simeq9.84519\units{m^2\,s^{-3}}$ is smaller than
$p_{A,0}(0)\simeq26.68598\units{m^2\,s^{-3}}$.
The defect changes sign near $z=0.00931273\units m$ and remains
positive thereafter, while its full-layer integral is zero.
These signs refer to the accepted budget cutoff; retaining the
original molecular cutoff would give a different imbalance.

Implementation should evaluate $K$ from the common matching trace,
construct $D_A$ with the calibrated coefficients, and determine
the cutoff from reference production at that same $K$.
Actual production is used in \cref{eq:accepted-k-transport} or
\cref{eq:accepted-curved-k}. Initialize a flat reference with
\cref{eq:accepted-k-profile}, retain zero modeled ground $k$ flux,
and couple the matching value and conservative flux to outer CLS.
Do not reintroduce the EP redistribution while using this profile.
Quadrature should resolve each active viscosity, forcing, and
dissipation breakpoint, and the full assembled reference residual
must be checked with the actual interface and pressure operators.

The derivation establishes continuum compatibility with the
standard outer CLS equilibrium on flat, undisturbed terrain and
specifies a conservative curved extension. It does not require
the outer flow over a hill to remain logarithmic. The positive
flat reference is also not a positivity or stability theorem for
the evolving equations: a sink controlled by sampled $K$ can
remain positive as local $k_w$ decreases, and the constitutive
limitations of the unchanged outer CLS model identified in
\cref{sec:stability} remain. Any active robustness modification
must be assessed against this reference and declared separately.
